\chapter{Topological, convex, and measure-theoretic tools}
\index{net}\index{upper envelope}\index{Radon measure}
\label{chap:topological-tools}

The spaces of metrics used later are seldom finite-dimensional and are not
always metrizable.  Suprema of plurisubharmonic functions need not be upper
semicontinuous, weak convergence of measures is topological rather than
pointwise, and several natural functionals take the value $-\infty$ or
$+\infty$.  These elementary facts account for a surprising number of gaps in
informal arguments.  The purpose of this chapter is to isolate the tools that
will be used repeatedly and to fix safe conventions once and for all.

Three failures recur, and it is worth naming them before any definition is
given.

\emph{First, sequences do not see the topology.}  A Berkovich analytification
of a proper variety is compact and Hausdorff, but it is metrizable only under
countability hypotheses on the ground field and the variety.  A compact space
that is not metrizable can contain a sequence with no convergent subsequence;
\cref{ex:net-no-subsequence} exhibits one explicitly.  Any argument of the
shape ``by compactness, pass to a convergent subsequence'' is therefore not a
proof until metrizability, first countability, or a Choquet reduction has been
established.  We consistently argue with nets, and we mark each place where a
sequence is legitimately recovered.

\emph{Second, suprema destroy upper semicontinuity.}  Plurisubharmonic
functions are upper semicontinuous by definition, and every envelope
construction in this book forms the supremum of a family of them.  That
supremum is in general not upper semicontinuous, as
\cref{ex:sup-not-usc,ex:sup-not-usc-log} show, and the correct object is its
upper semicontinuous regularization.  The whole content of the envelope
theorem of \cref{part:bj} is that this regularization stays inside the psh
class; the regularization itself is a purely topological operation, developed
here.

\emph{Third, the natural functionals are extended-real.}  Potentials take the
value $-\infty$, energies of unbounded potentials may be $-\infty$, dual
energies may be $+\infty$, and differences of two such quantities are
routinely written down in the literature.  \cref{sec:extended-arithmetic}
fixes a ledger for these operations, justifies each entry, and shows by
explicit example what fails if the entry is dropped.

The chapter depends only on elementary topology and measure theory; no
Berkovich geometry is used in the proofs, although the examples are chosen so
that their later incarnations are visible.  The
capacity language in \cref{sec:capacity-preview} is a preview: concrete
Monge--Amp\`ere capacities will be constructed in
\cref{part:bfj,part:bj}.  Likewise,
the concave transforms in \cref{sec:dh-preview} anticipate filtrations and
test curves without presupposing them.

\subsection*{Relation to the appendices}

Two appendices carry the same material at a different level of generality, and
the division of labour is deliberate.
\Cref{app:nets} is a reference sheet on directed sets: it fixes the
order-preserving tail convention for subnets, proves the cluster-point and
compactness equivalences for that convention, and records the extended-real
identities in their most general form.
\Cref{app:convexity} is a reference sheet on convex analysis: closed convex
functions on a finite-dimensional vector space, Legendre--Fenchel duality,
mixed volumes, and Alexandrov Monge--Amp\`ere measures.

The present chapter is not a summary of either.  It states, in the exact form
in which they will be quoted, the finitely many statements that the main
narrative uses; it proves those statements; and it supplies the worked
examples and counterexamples that show why the hypotheses cannot be weakened.
When a variant is needed only once, or only in a technical step, the
statement is quoted from the appendix rather than reproved.  In particular we
use the Willard subnet convention of \Cref{def:appE-willard-subnet}
throughout, and we use the finite-dimensional biconjugation theorem
\Cref{thm:appD-fenchel-moreau} as the one imported convex-analytic input of
\cref{sec:radon-weak}.

\subsection*{Roadmap}

\Cref{sec:nets-compactness} develops Moore--Smith convergence: nets, subnets,
the net characterizations of closure, continuity, and compactness, the net
description of semicontinuity, monotone nets, and Dini's theorem in the two
forms that occur later.
\Cref{sec:usc-envelopes} treats upper semicontinuous regularization, proves
the identities that the envelope calculus uses without comment, and proves the
Choquet selection lemma together with a sharp analysis of the countability
hypothesis.
\Cref{sec:radon-weak} states the Riesz representation interface precisely,
develops weak convergence of Radon measures, and proves the finite-dimensional
energy duality and Legendre transform statements that the variational chapters
imitate in infinite dimensions.
\Cref{sec:capacity-preview} sets up capacities, quasicontinuity, and polar
sets axiomatically, at exactly the level of generality that \cref{part:bfj}
and \cref{part:bj} will import.
\Cref{sec:dh-preview} treats concave functions on convex bodies and their
Duistermaat--Heckman distributions.
\Cref{sec:extended-arithmetic} contains the extended-real ledger.
The closing section collects the dictionary and the sources.

\section{Nets and compactness}
\label{sec:nets-compactness}

\begin{definition}
\label{def:directed-net}
A \emph{directed set} is a nonempty preordered set $(A,\leq)$ such that for any
$a_1,a_2\in A$ there exists $a_3\in A$ with $a_1,a_2\leq a_3$.  A \emph{net}
in a set $X$ indexed by $A$ is a map $A\to X$, written $(x_a)_{a\in A}$.
If $X$ is a topological space, the net converges to $x\in X$ if for every
neighborhood $U$ of $x$ there exists $a_0\in A$ such that $x_a\in U$ whenever
$a\geq a_0$.
\end{definition}

Sequences are nets indexed by $\mathbf N$, but they do not detect the topology
of an arbitrary space.  We will therefore state compactness and continuity in
terms of nets unless a metrizability or countability argument has already
reduced the discussion to sequences.

Thus antisymmetry is not required; this agrees with
\Cref{def:appE-directed}, and some of the index sets occurring in this book
are naturally preordered rather than ordered.  Two pieces of vocabulary
shorten every later statement.  A property $P(a)$ holds
\emph{eventually} if there is $a_0$ with $P(a)$ for all $a\geq a_0$, and it
holds \emph{frequently} if for every $a_0$ there is $a\geq a_0$ with $P(a)$.
Thus $x_a\to x$ means that $x_a$ is eventually in every neighborhood of $x$,
and $x$ is a \emph{cluster point} of $(x_a)$ if $x_a$ is frequently in every
neighborhood of $x$.  Negating ``eventually'' gives ``frequently not'', a
rewriting used repeatedly and silently.

The directed sets that actually occur in this book are worth listing once,
because their directedness is exactly the property that makes the later
arguments work.

\begin{example}[The directed sets used in this book]
\label{ex:directed-sets-of-the-book}
Each of the following is directed, for the stated reason.
\begin{enumerate}
\item The natural numbers with their usual order.  Nets indexed by $\mathbf N$ are
  sequences.
\item The neighborhoods of a point $x$ of a topological space, ordered by reverse
  inclusion.  Two neighborhoods have their intersection as a common upper
  bound.  This is the index set of the neighborhood net of
  \Cref{ex:appE-neighborhood-net}, and it is the device that replaces
  ``choose a sequence tending to $x$'' in a space with no countable
  neighborhood basis.
\item The finite subsets of a set $S$, ordered by inclusion.  Two finite subsets
  have their union as an upper bound.  This index set produces the standard
  counterexamples of \cref{sec:extended-arithmetic} and of
  \Cref{ex:appE-measurable-net-failure}.
\item The models of a fixed variety over a valuation ring, ordered by domination
  as in \cref{chap:models}.  Two models are dominated by a common third one,
  obtained by normalizing the closure of the diagonal in their fibre product.
  Model functions and model metrics are indexed by this directed set, and
  \cref{chap:bfj-psh} takes limits along it.
\item A family $\mathcal F$ of functions that is stable under finite maxima,
  ordered pointwise.  The maximum of two members is a common upper bound.
  Every envelope in this book is the supremum over such a family, and
  directedness is what converts a supremum into an increasing limit.
\item The pairs $(a,\varepsilon)\in A\times\mathbf R_{>0}$, where $A$ is directed,
  ordered by $a\leq a'$ and $\varepsilon\geq\varepsilon'$.  This is the index
  set of every diagonal argument in which one index is refined while a
  tolerance is decreased.
\end{enumerate}
Directedness fails, by contrast, for the family of \emph{all} subsets of $S$
of cardinality exactly one, since $\{s\}$ and $\{s'\}$ have no common upper
bound of cardinality one; that family is not used as an index set.  Note
that the compact subsets of any space are directed by inclusion, since
$K_1\cup K_2$ is a compact upper bound.  This is the exhaustion order used
in \cref{sec:radon-weak}; reverse inclusion instead records shrinking
intersections and is not an exhaustion order.
\end{example}

\begin{definition}
\label{def:subnet}
Let $(x_a)_{a\in A}$ be a net.  A \emph{subnet} is a net
$(x_{\sigma(b)})_{b\in B}$, where $B$ is directed and
$\sigma:B\to A$ is order preserving and cofinal.  For an order-preserving
map, cofinality is equivalent to the following tail condition: for every
$a_0\in A$ there exists $b_0\in B$ such that
\[
                   b\geq b_0\quad\Longrightarrow\quad
                   \sigma(b)\geq a_0.
\]
No injectivity of $\sigma$ is required.
\end{definition}

The eventual-cofinality condition is essential.  Merely choosing infinitely
many terms of a net can discard an entire tail and does not give a subnet in
the topological sense.

\begin{remark}[Which subnet convention is in force]
\label{rem:subnet-conventions}
Terminology for subnets varies between topology texts, so the displayed
formula, rather than the name, controls every use in this book.  Our map
$\sigma$ is order preserving and is eventually in every tail of $A$; this is
exactly the convention of \Cref{def:appE-willard-subnet,def:appE-ms-subnet}.
With this convention \Cref{thm:appE-cluster-subnet} characterizes cluster
points by convergent subnets.  What is never permissible is to replace a
subnet by a subsequence: a sequence of selected terms can be cofinal only if
there is a tail-cofinal map $\mathbf N\to A$, and
\cref{ex:no-cofinal-sequence} shows that such a map need not exist.
\end{remark}

\begin{example}[A directed set with no cofinal sequence]
\label{ex:no-cofinal-sequence}
Let $S$ be an uncountable set and let $A$ be the collection of finite subsets
of $S$, directed by inclusion as in
\Cref{ex:directed-sets-of-the-book}.  Suppose $\sigma:\mathbf N\to A$ were
cofinal, so that every $F\in A$ satisfies $F\subseteq\sigma(n)$ for some $n$.
Then $S=\bigcup_{n}\sigma(n)$ is a countable union of finite sets, hence
countable, contradicting the choice of $S$.  Consequently a net indexed by
$A$ has no subnet indexed by $\mathbf N$, and every statement about such a net
must be proved for nets.  This directed set is not exotic: it is the index set
of the finite-subset net in \Cref{ex:appE-measurable-net-failure} and, in
\cref{part:bj}, of the family of finite maxima of Fubini--Study functions
attached to a possibly uncountable family of sections.
\end{example}

\begin{proposition}\label{prop:net-criterion}
Let $X$ and $Y$ be topological spaces.
\begin{enumerate}
\item A subset $F\subset X$ is closed if and only if every convergent net in
  $F$ has its limit in $F$.
\item A map $f:X\to Y$ is continuous if and only if $x_a\to x$ implies
  $f(x_a)\to f(x)$ for every net $(x_a)$.
\item The space $X$ is compact if and only if every net in $X$ has a
  convergent subnet.
\item If $X$ is Hausdorff, a convergent net has a unique limit.
\end{enumerate}
\par
\end{proposition}

\begin{proof}
We prove all four assertions; the neighborhood net of
\Cref{ex:appE-neighborhood-net} is the only construction used in the first
two.

\emph{(1).}  Suppose $F$ is closed, let $(x_a)$ be a net in $F$, and let
$x_a\to x$.  If $x\notin F$, then $X\setminus F$ is a neighborhood of $x$, so
the net is eventually in $X\setminus F$; but every term lies in $F$, a
contradiction.  Conversely, suppose $F$ is not closed and choose
$x\in\overline F\setminus F$.  Direct the neighborhoods of $x$ by reverse
inclusion; since $x\in\overline F$, every neighborhood $U$ meets $F$, so we
may choose $x_U\in F\cap U$.  Given a neighborhood $V$ of $x$, every
$U\preceq$-larger than $V$ satisfies $U\subseteq V$ and hence $x_U\in V$.
Thus $x_U\to x$, and this convergent net in $F$ has its limit outside $F$.

\emph{(2).}  Let $f$ be continuous, let $x_a\to x$, and let $V$ be a
neighborhood of $f(x)$.  Then $f^{-1}(V)$ is a neighborhood of $x$, so $x_a$
is eventually in it and $f(x_a)$ is eventually in $V$.  Conversely, suppose
$f$ is not continuous at $x$: some neighborhood $V$ of $f(x)$ has the property
that no neighborhood $U$ of $x$ satisfies $f(U)\subseteq V$.  Equivalently,
every neighborhood $U$ of $x$ contains a point $x_U$ with
$f(x_U)\notin V$.  The resulting neighborhood net converges to $x$ by the
argument in (1), while $f(x_U)$ never enters $V$ and so does not converge to
$f(x)$.

\emph{(3).}  Suppose first that $X$ is compact and let $(x_a)_{a\in A}$ be a
net.  Direct the tails of the net together with their closures:
\[
                  F_a=\overline{\{x_b:b\geq a\}}.
\]
Directedness gives the finite-intersection property, so compactness yields
$x\in\bigcap_aF_a$.  For every pair $(a,U)$, with $U$ a neighborhood of
$x$, there is some $b\geq a$ such that $x_b\in U$.  Let $B$ consist of all
triples $(a,U,b)$ with these properties, and declare
\[
 (a,U,b)\leq(a',U',b')
 \quad\Longleftrightarrow\quad
 a\leq a',\quad U'\subseteq U,\quad b\leq b'.
\]
The cluster-point property shows that $B$ is directed: first dominate the
two $a$- and $b$-indices, intersect the two neighborhoods, and then choose a
later term in that intersection.  The map
$\sigma(a,U,b)=b$ is order preserving and cofinal, and the associated subnet
converges to $x$.

We spell out the last three assertions of that paragraph, since they are the
only delicate point.  The set $B$ is nonempty over each pair $(a,U)$ precisely
because $x\in F_a$ for every $a$, so the tail $\{x_b:b\geq a\}$ meets $U$.
Given $(a_1,U_1,b_1)$ and $(a_2,U_2,b_2)$, choose $a_3\geq a_1,a_2,b_1,b_2$ in
$A$, put $U_3=U_1\cap U_2$, and use $x\in F_{a_3}$ to find $b_3\geq a_3$ with
$x_{b_3}\in U_3$; then $(a_3,U_3,b_3)$ dominates both triples.  The map
$\sigma(a,U,b)=b$ is order preserving by construction.  It is cofinal in the
tail sense: given $a_0\in A$, pick any $U_0\ni x$ and $b_0\geq a_0$ with
$x_{b_0}\in U_0$; every triple $\geq(a_0,U_0,b_0)$ has last entry
$\geq b_0\geq a_0$.  Finally the subnet converges to $x$: given a neighborhood
$V$ of $x$, choose $b_V$ with $x_{b_V}\in V$ and consider triples
$\geq(b_V,V,b_V)$; their neighborhood entry $U$ satisfies $U\subseteq V$, and
the defining condition $x_b\in U$ gives $x_b\in V$.

Conversely, if an open cover $\mathcal U$ has no finite
subcover, direct its finite subfamilies by inclusion and choose
\[
                  x_{\mathcal V}\notin\bigcup_{U\in\mathcal V}U.
\]
No subnet can converge: a limit belongs to some $U\in\mathcal U$, whereas
eventual cofinality forces the subnet eventually to have indices containing
$U$ and hence to lie outside $U$.

\emph{(4).}  Let $X$ be Hausdorff and suppose $x_a\to x$ and $x_a\to y$ with
$x\neq y$.  Choose disjoint neighborhoods $U\ni x$ and $V\ni y$.  The net is
eventually in $U$ and eventually in $V$; directedness produces one index
beyond both thresholds, and the corresponding term lies in $U\cap
V=\varnothing$.
\end{proof}

\begin{remark}
\label{rem:nets-safe-default}
On a first-countable space, limits are detected by sequences, and on a metric
space compactness is equivalent to sequential compactness.  Berkovich spaces
that occur in this book are compact Hausdorff under properness hypotheses,
but they need not come with a preferred metric.  The net formulation is
therefore the safe default.
\end{remark}

The following example is the reason the previous sentence is not a stylistic
preference.  It shows that a compact Hausdorff space can carry a sequence with
no convergent subsequence, so that the phrase ``extract a convergent
subsequence'' has no content there, while
\Cref{prop:net-criterion} still supplies a convergent subnet.

\begin{example}[A sequence in a compact space with no convergent subsequence]
\label{ex:net-no-subsequence}
Let
\[
 X=\{0,1\}^{[0,1]}
 =\bigl\{\text{all maps }[0,1]\to\{0,1\}\bigr\},
\]
with the product topology, where $\{0,1\}$ is discrete.  By the Tychonoff
theorem \cite[Chapter~5]{Kel75}, $X$ is compact; it is Hausdorff because two
distinct maps differ at some $t$, and the corresponding coordinate projections
separate them.  Convergence in $X$ is exactly pointwise convergence:
a basic neighborhood constrains finitely many coordinates.

For $t\in[0,1)$ let
\[
 t=\sum_{n\geq1}\beta_n(t)2^{-n},\qquad\beta_n(t)\in\{0,1\},
\]
be the binary expansion, chosen to be the one that is eventually $0$ when $t$
is a dyadic rational.  Define $\beta_n(1)=0$ for every $n$; then $\beta_n$
is a well-defined map $[0,1]\to\{0,1\}$.  Set $x_n=\beta_n\in X$.

We claim that no subsequence of $(x_n)$ converges.  Let
$n_1<n_2<n_3<\cdots$ be any strictly increasing sequence of indices and put
\[
 t_0=\sum_{k\geq1}2^{-n_{2k}}\in[0,1].
\]
Since the exponents $n_2<n_4<n_6<\cdots$ are strictly increasing, this series
is a binary expansion that is not eventually $1$, so it is the chosen
expansion of $t_0$ and
\[
 \beta_{n_{2k}}(t_0)=1,\qquad \beta_{n_{2k-1}}(t_0)=0
 \qquad(k\geq1).
\]
Thus the coordinate of $x_{n_j}$ at $t_0$ alternates between $0$ and $1$ and
does not converge.  Since convergence in $X$ is coordinatewise, the
subsequence $(x_{n_j})$ does not converge.

By \Cref{prop:net-criterion}\,(3) the sequence $(x_n)$ nevertheless has a
convergent subnet; concretely, $\{0,1\}^{[0,1]}$ being compact, the closed
tails $\overline{\{x_m:m\geq n\}}$ have a common point $x$, and the
construction in the proof produces a subnet converging to $x$.  No such subnet
is indexed by $\mathbf N$ in an order-preserving cofinal way, by the argument
just given.

The example is not pathological from the point of view of this book.  The
space $X$ is the space of $\{0,1\}$-valued functions on an uncountable set;
spaces of the same nature arise as spaces of semivaluations, and
$\mathcal P(X)$ with its weak topology, studied in \cref{sec:radon-weak}, is a
compact convex set that is metrizable only when $C(X)$ is separable.
\end{example}

Semicontinuity is the property that the book tests most often, and it too has
a clean net formulation.  Recall that $f:X\to\Rbar$ is \emph{upper
semicontinuous}, abbreviated usc, if $\{f<c\}$ is open for every
$c\in\mathbf R$, and \emph{lower semicontinuous}, abbreviated lsc, if $-f$ is
usc.

\begin{proposition}\label{prop:net-semicontinuity}
Let $X$ be a topological space and $f:X\to\Rbar$.  Then $f$ is upper
semicontinuous if and only if
\begin{equation}
\label{eq:net-usc-test}
 \limsup_a f(x_a)\leq f(x)
\end{equation}
for every net $(x_a)$ in $X$ with $x_a\to x$.  Dually, $f$ is lower
semicontinuous if and only if $\liminf_a f(x_a)\geq f(x)$ for every such net.
\end{proposition}

\begin{proof}
Assume $f$ is usc and let $x_a\to x$.  If $f(x)=+\infty$ there is nothing to
prove, so let $c>f(x)$ be real.  Then $\{f<c\}$ is an open neighborhood of
$x$, so $f(x_a)<c$ eventually, whence $\limsup_af(x_a)\leq c$.  Letting
$c\downarrow f(x)$ gives \eqref{eq:net-usc-test}; if $f(x)=-\infty$ the same
argument with $c\downarrow-\infty$ gives $\limsup_af(x_a)=-\infty$.

Conversely, suppose $f$ is not usc.  Then $\{f<c\}$ fails to be open for some
real $c$, so there is $x$ with $f(x)<c$ such that every neighborhood $U$ of
$x$ contains a point $x_U$ with $f(x_U)\geq c$.  The neighborhood net
$(x_U)$ converges to $x$ by the proof of \Cref{prop:net-criterion}\,(1),
while
\[
 \limsup_U f(x_U)\geq c>f(x),
\]
contradicting \eqref{eq:net-usc-test}.  The lower semicontinuous case follows
by replacing $f$ with $-f$ and using
$\liminf_a(-r_a)=-\limsup_a r_a$, which is immediate from the definitions in
\cref{sec:extended-arithmetic}.
\end{proof}

\begin{corollary}\label{cor:usc-attains-max}
Let $K$ be a nonempty compact space and let
$f:K\to\mathbf R\cup\{-\infty\}$ be upper semicontinuous.  Then $f$ attains
its supremum; in particular $f$ is bounded above.
\end{corollary}

\begin{proof}
Put $M=\sup_Kf\in\mathbf R\cup\{\pm\infty\}$.  If $M=-\infty$ then
$f\equiv-\infty$ and every point is a maximizer, so assume $M>-\infty$ and let
$A=\{m\in\mathbf R:m<M\}$, directed by the usual order and nonempty.  For
$m\in A$ choose $x_m\in K$ with $f(x_m)>m$, which is possible by the
definition of a supremum.  For $m\geq m_0$ we get $f(x_m)>m_0$, so
\[
 \liminf_{m}f(x_m)\geq\sup_{m_0\in A}m_0=M,
\]
while $\limsup_mf(x_m)\leq M$ trivially; hence $f(x_m)\to M$.  By
\Cref{prop:net-criterion}\,(3) there is a subnet $x_{\sigma(b)}\to x\in K$,
and by \Cref{lem:appE-tail-subnet} the corresponding subnet of values still
converges to $M$.  \Cref{prop:net-semicontinuity} gives
\[
 f(x)\geq\limsup_bf(x_{\sigma(b)})=M .
\]
Since $f\leq M$ everywhere, $f(x)=M$.  As $f$ does not take the value
$+\infty$, this forces $M\in\mathbf R$, so $f$ is bounded above.
\end{proof}

\subsection{Monotone nets of functions}

Let $S$ be a set.  A net $(f_a)$ of extended-real functions on $S$ is
\emph{increasing} if $a\leq b$ implies $f_a\leq f_b$ pointwise, and
\emph{decreasing} if $a\leq b$ implies $f_a\geq f_b$.  We then write
$f_a\nearrow f$ or $f_a\searrow f$ for the pointwise supremum or infimum.

A monotone net of extended-real numbers always converges in $\Rbar$, to its
supremum or infimum; this is the one convergence statement that never needs a
compactness hypothesis, and we record it because it is used at almost every
occurrence of a monotone family.

\begin{lemma}\label{lem:monotone-net-converges}
Let $(r_a)_{a\in A}$ be an increasing net in $\Rbar$ and let $r=\sup_ar_a$.
Then $r_a\to r$ in the order topology of $\Rbar$.  The analogous statement
holds for decreasing nets and infima.  Consequently, for a monotone net
$(f_a)$ of extended-real functions on a set $S$, the pointwise limit exists at
every point and equals the pointwise supremum, respectively infimum.
\end{lemma}

\begin{proof}
A neighborhood basis of $r$ in $\Rbar$ consists of the sets $(c,d)$ with
$c<r<d$ when $r$ is finite, of the sets $(c,+\infty]$ with $c<r$ when
$r=+\infty$, and of the sets $[-\infty,d)$ with $d>r$ when $r=-\infty$.  In
all three cases it suffices to show that for $c<r$ the net is eventually in
$(c,+\infty]$ and that for $d>r$ it is eventually in $[-\infty,d)$.  The
second statement holds for every index, since $r_a\leq r<d$.  For the first,
$c<r=\sup_ar_a$ produces $a_0$ with $r_{a_0}>c$, and monotonicity gives
$r_a\geq r_{a_0}>c$ for $a\geq a_0$.  The decreasing case follows by applying
this to $(-r_a)$, and the pointwise statement is the numerical statement at
each point of $S$.
\end{proof}

\begin{proposition}\label{prop:dini-nets}
Let $K$ be compact, and let $(f_a)$ be a decreasing net of continuous
real-valued functions on $K$.  If $f_a\searrow f$ and $f$ is continuous,
then $f_a\to f$ uniformly.
\end{proposition}

\begin{proof}
Fix $\varepsilon>0$.  The sets
\[
                     U_a=\{x\in K:f_a(x)<f(x)+\varepsilon\}
\]
are open, increasing with $a$, and cover $K$.  A finite subcover is dominated
by one index $a_0$.  Hence $0\leq f_a-f<\varepsilon$ on $K$ for all
$a\geq a_0$.
\end{proof}

Three points in this proof deserve to be isolated, because the form of Dini's
theorem actually used later weakens the hypotheses in exactly these places.
The sets $U_a$ are open because $f_a-f$ is continuous; they increase because
the net decreases; and they cover $K$ because $\inf_af_a=f$ pointwise.  The
finite subcover is dominated by a single index because the family $(U_a)$ is
directed, so that a finite union $U_{a_1}\cup\cdots\cup U_{a_k}$ is contained
in $U_{a_0}$ for any $a_0\geq a_1,\dots,a_k$.  Openness of $U_a$ survives if
$f_a$ is merely usc and $f$ is merely lsc, since then $f_a-f$ is usc where the
difference is defined.  This gives the version we shall quote.

\begin{proposition}
\label{prop:dini-usc}
Let $K$ be compact, let $(f_a)$ be a decreasing net of upper semicontinuous
functions $K\to\mathbf R\cup\{-\infty\}$, and let $f=\inf_af_a$.
\begin{enumerate}
\item If $f$ is lower semicontinuous and real valued, then $f_a\to f$
  uniformly on $K$, and in particular $f$ is continuous and $f_a$ is
  eventually real valued and bounded.
\item Without further hypotheses, and for every real $c$,
\[
 f<c\ \text{ on }K
 \quad\Longrightarrow\quad
 f_a<c\ \text{ on }K\ \text{ eventually}.
\]
\end{enumerate}
\par
\end{proposition}

\begin{proof}
\emph{(2).}  The sets $V_a=\{f_a<c\}$ are open because $f_a$ is usc, and they
increase with $a$ because the net decreases.  If $f<c$ everywhere, then for
each $x$ there is $a$ with $f_a(x)<c$, so the $V_a$ cover $K$.  Compactness
gives a finite subcover, and directedness gives one index $a_0$ with
$V_{a_0}=K$.  Monotonicity gives $V_a=K$ for $a\geq a_0$.

\emph{(1).}  Fix $\varepsilon>0$ and set
\[
 U_a=\{x\in K:f_a(x)<f(x)+\varepsilon\}.
\]
Since $f$ is lsc and real valued and $f_a$ is usc, the function $f_a-f$ is
usc, so $U_a$ is open; the sets increase with $a$; and they cover $K$ because
$\inf_af_a=f$ pointwise with $f(x)$ finite.  As before, compactness and
directedness give $a_0$ with $U_{a_0}=K$, and then
\[
 0\leq f_a-f<\varepsilon\qquad\text{on }K,\ a\geq a_0 .
\]
In particular $f_{a_0}$ is real valued, and $f$, being a uniform limit of the
usc functions $f_a$, is usc; being lsc by hypothesis, it is continuous.
\end{proof}

\begin{remark}
\label{rem:dini-hypotheses-sharp}
Both hypotheses in \Cref{prop:dini-usc}\,(1) are needed.  Compactness cannot
be dropped: on $\mathbf R$ the sequence $f_n(x)=\max(0,1-|x-n|)$ decreases to
$0$ after replacing $f_n$ by $g_n=\sup_{m\geq n}f_m$, and $g_n\searrow0$
pointwise while $\sup|g_n|=1$ for every $n$.  Semicontinuity of the limit
cannot be dropped: on $K=[0,1]$ the functions
$f_n(x)=\min(1,nx)$ form an increasing sequence of continuous functions, so
$h_n=1-f_n$ decreases to $h=\mathbf 1_{\{0\}}$, which is usc but not lsc, and
$\sup_K(h_n-h)=1$ for every $n$.  In the second example, part (2) of the
proposition still holds, since $h<c$ on $K$ forces $c>1\geq h_n$; it is
therefore part (2) that the book uses whenever the limit is not known in
advance to be continuous.  In the first example part (2) fails as well, which
shows that compactness is used in both parts and not merely in the uniform
estimate.
\end{remark}

\begin{remark}[Warning]
\label{rem:monotone-net-integration-warning}
The measure-theoretic monotone convergence theorem does not hold for
arbitrary nets of measurable functions.  Let $X=[0,1]$ with Lebesgue measure,
direct the finite subsets $F\subset X$ by inclusion, and put
$f_F=\mathbf 1_F$.  Then $f_F\nearrow1$ pointwise, while
$\int f_F=0$ for every $F$.  Whenever an argument integrates a monotone net,
it must use additional regularity or first reduce to a sequence.
\end{remark}

The following regular version is often sufficient in pluripotential theory.

\begin{proposition}\label{prop:lsc-net-monotone}
Let $K$ be compact Hausdorff, let $\mu$ be a finite Radon measure, and let
$(f_a)$ be an increasing net of nonnegative lower semicontinuous functions.
Then
\[
             \int_K\sup_a f_a\,\rd \mu
             =\sup_a\int_K f_a\,\rd \mu.
\]
\end{proposition}

\begin{proof}
Write $f=\sup_a f_a$, which is lower semicontinuous.  For a nonnegative lower
semicontinuous function on a compact Hausdorff space,
\[
 \int_K f\,\rd \mu
 =\sup\left\{\int_K h\,\rd \mu:h\in C(K),\ 0\leq h\leq f\right\}.
\]
Fix such an $h$ and $\varepsilon>0$.  The open sets
$\{f_a>h-\varepsilon\}$ cover $K$.  Compactness and directedness give one
index $a$ for which $f_a>h-\varepsilon$ everywhere.  Thus
\[
 \sup_a\int_K f_a\,\rd \mu
 \geq\int_Kh\,\rd \mu-\varepsilon\mu(K).
\]
Taking the supremum over $h$ and then letting $\varepsilon\downarrow0$ proves
the nontrivial inequality.
\end{proof}

Two comments make the proof reusable.  First, the inequality $\leq$ is the
easy one: each $f_a\leq f$ gives $\int f_a\leq\int f$, so the supremum on the
right is at most the left-hand side.  Second, the displayed characterization
of $\int f$ by continuous minorants is the Radon regularity statement proved
in detail as \eqref{eq:appE-lsc-integral-minorants}; it is where the
hypothesis that $\mu$ be Radon, rather than merely Borel, is used.  The
covering step then uses only that $\{f_a>h-\varepsilon\}$ is open, which is
lower semicontinuity of $f_a$, and that the family is directed.

The form in which we shall usually quote this statement is the dual one, for
decreasing nets of upper semicontinuous functions.  Since potentials are usc
and decreasing approximations are the standard regularization device, this is
the version that appears in the construction of Monge--Amp\`ere measures and
energies.

\begin{corollary}\label{cor:usc-net-monotone-integral}
Let $K$ be compact Hausdorff, let $\mu$ be a finite Radon measure on $K$, and
let $(f_a)$ be a decreasing net of upper semicontinuous functions
$K\to\mathbf R\cup\{-\infty\}$ that are uniformly bounded above by a constant
$C$.  Then
\[
 \int_K\inf_af_a\,\rd\mu=\inf_a\int_Kf_a\,\rd\mu,
\]
both sides being well defined in $[-\infty,+\infty)$ and possibly equal to
$-\infty$.
\end{corollary}

\begin{proof}
The functions $g_a=C-f_a$ take values in $[0,+\infty]$, are lower
semicontinuous, and form an increasing net with
$\sup_ag_a=C-\inf_af_a$.  \Cref{prop:lsc-net-monotone} gives
\[
 \int_K\bigl(C-\inf_af_a\bigr)\rd\mu=\sup_a\int_K(C-f_a)\,\rd\mu .
\]
Each $f_a$ is bounded above by $C$, so $\int_Kf_a\,\rd\mu$ is well defined in
$[-\infty,+\infty)$ and
$\int_K(C-f_a)\,\rd\mu=C\mu(K)-\int_Kf_a\,\rd\mu$ as an identity in
$[0,+\infty]$; the same applies to $\inf_af_a$.  Subtracting the finite
number $C\mu(K)$ and changing signs, which is legitimate because only one
infinite value can occur and it occurs on the same side of each equation,
yields the assertion.  This is the statement recorded as
\Cref{cor:appE-usc-monotone}.
\end{proof}

The next example shows that a decreasing net of usc functions can converge to
its infimum in a way that no sequence extracted from the net reproduces.  It
is the exact counterpart, for infima, of the fact that suprema require
regularization, and it explains why the monotone approximation theorems of
\cref{part:bfj,part:bj} are stated for nets and why their sequential
corollaries always carry a countability hypothesis.

\begin{example}[A decreasing net of usc functions with no cofinal sequence]
\label{ex:decreasing-net-not-sequential}
Let $S$ be an uncountable set and let
\[
 X=S\cup\{\infty\}
\]
be its one-point compactification: every subset of $S$ is open, and the
neighborhoods of $\infty$ are the complements of the finite subsets of $S$.
Then $X$ is compact, since a cover of $X$ contains a member $U\ni\infty$ whose
complement is finite, and Hausdorff, since two points of $S$ are separated by
singletons and a point $s\in S$ is separated from $\infty$ by $\{s\}$ and
$X\setminus\{s\}$.  It is not first countable at $\infty$: a countable family
of neighborhoods of $\infty$ has complements with countable union, so any
$s\in S$ outside that union gives a neighborhood $X\setminus\{s\}$ containing
no member of the family.

Direct the finite subsets $F\subset S$ by inclusion and put
\[
 f_F=\mathbf 1_{X\setminus F}.
\]
Each $F$ is open, so $X\setminus F$ is closed and $f_F$ is upper
semicontinuous: for $c\leq0$ the set $\{f_F<c\}$ is empty, for $0<c\leq1$ it
is $F$, which is open, and for $c>1$ it is $X$.  The net is decreasing,
because $F\subseteq F'$ gives $f_{F'}\leq f_F$.  Its pointwise infimum is
\[
 f=\inf_Ff_F=\mathbf 1_{\{\infty\}},
\]
since $f_{\{s\}}(s)=0$ for $s\in S$ while $f_F(\infty)=1$ for every finite
$F\subset S$.  The limit $f$ is usc, in agreement with
\Cref{lem:usc-inf}.

Now let $(F_j)_{j\geq1}$ be any sequence of indices.  The union
$\bigcup_jF_j$ is countable, so there exists $s\in S$ with $s\notin F_j$ for
all $j$, and then
\[
 \inf_jf_{F_j}(s)=1\neq0=f(s).
\]
Thus no sequence extracted from the net decreases to $f$, and by
\Cref{ex:no-cofinal-sequence} the index set admits no cofinal sequence at all.
Note also that $f_F\searrow f$ while $\sup_X(f_F-f)=1$ for every $F$, so the
uniform conclusion of \Cref{prop:dini-usc}\,(1) genuinely fails here; indeed
$f=\mathbf 1_{\{\infty\}}$ is not lower semicontinuous, because
$\{f>1/2\}=\{\infty\}$ is not open.
\end{example}

\begin{remark}
\label{rem:net-vs-sequence-checklist}
\Cref{ex:net-no-subsequence,ex:no-cofinal-sequence,ex:decreasing-net-not-sequential}
form the standard test battery for the phrase ``pass to a subsequence''.  A
sequential reduction in this book is always justified by exactly one of three
mechanisms: metrizability or second countability of the underlying space, as
in \Cref{prop:measure-compactness}; the Choquet selection lemma
\Cref{lem:choquet-selection}, which replaces a family by a countable
subfamily with the same regularized supremum; or an explicit cofinal sequence
in the index set, as when models are indexed by a countable system of blowups.
No other mechanism is used.
\end{remark}

\section{Upper semicontinuous envelopes}
\label{sec:usc-envelopes}

Let $X$ be a topological space and
$f:X\to\overline{\mathbf R}:=\mathbf R\cup\{-\infty,+\infty\}$.  Its
\emph{upper semicontinuous regularization} is
\begin{equation}
\label{eq:usc-regularization}
 f^*(x)=\inf_{U\ni x}\ \sup_{y\in U}f(y),
\end{equation}
where $U$ runs over the neighborhoods of $x$.  The infimum is over a directed
family, decreasing as $U$ shrinks, so it is a monotone limit in the sense of
\Cref{lem:monotone-net-converges}; in particular it always exists in $\Rbar$
and no measurability or finiteness assumption is needed.  It suffices to let
$U$ run over any neighborhood basis of $x$.

\begin{proposition}
\label{prop:usc-majorant}
The function $f^*$ is upper semicontinuous, satisfies $f\leq f^*$, and is the
smallest upper semicontinuous function majorizing $f$.  In particular,
$f=f^*$ if and only if $f$ is upper semicontinuous.
\end{proposition}

\begin{proof}
For $c\in\mathbf R$, the condition $f^*(x)<c$ means that
$\sup_U f<c$ for some neighborhood $U$ of $x$; it therefore persists on $U$.
Thus $\{f^*<c\}$ is open.  Since every neighborhood of $x$ contains $x$, we
have $f(x)\leq f^*(x)$.  If $g$ is upper semicontinuous and $f\leq g$, then
for every $c>g(x)$ there is a neighborhood $U$ of $x$ on which $g<c$.
Consequently $f^*(x)\leq c$, and $f^*(x)\leq g(x)$ after $c\downarrow g(x)$.

For the last sentence: if $f=f^*$ then $f$ is usc by the first assertion,
and if $f$ is usc then $f$ is an usc majorant of itself, so minimality gives
$f^*\leq f$, whence equality.
\end{proof}

The regularization has an equivalent description in terms of upper limits,
which is the form used whenever a potential is estimated along an approaching
net.

\begin{proposition}
\label{prop:usc-limsup-description}
Let $X$ be a topological space, $f:X\to\Rbar$, and $x\in X$.  Then
\begin{equation}
\label{eq:usc-limsup-description}
 f^*(x)=\max\Bigl\{f(x),\ \sup\bigl\{\limsup_a f(x_a):
 (x_a)\text{ a net in }X,\ x_a\to x\bigr\}\Bigr\},
\end{equation}
and the inner supremum is attained by the neighborhood--tolerance refinement
constructed below.  If $x$ is not isolated and $f$ is arbitrary,
\eqref{eq:usc-limsup-description} may be rewritten as
\[
 f^*(x)=\max\Bigl\{f(x),\ \limsup_{y\to x}f(y)\Bigr\},
 \qquad
 \limsup_{y\to x}f(y):=\inf_{U\ni x}\ \sup_{y\in U\setminus\{x\}}f(y).
\]
\end{proposition}

\begin{proof}
Write $L$ for the right-hand side of \eqref{eq:usc-limsup-description}.

We first prove $f^*(x)\geq L$.  Certainly $f^*(x)\geq f(x)$.  Let $x_a\to x$
be a net and let $U$ be a neighborhood of $x$.  Then $x_a\in U$ eventually, so
$\sup_{b\geq a_0}f(x_b)\leq\sup_Uf$ for suitable $a_0$, whence
$\limsup_af(x_a)\leq\sup_Uf$.  Taking the infimum over $U$ gives
$\limsup_af(x_a)\leq f^*(x)$.

For $f^*(x)\leq L$ we may assume $f^*(x)>-\infty$, the case
$f^*(x)=-\infty$ being trivial.  Then $\sup_Uf\geq f^*(x)>-\infty$ for every
neighborhood $U$ of $x$, so the choices below are possible.  Index a net by
the pairs $(U,\varepsilon)$, where $U$ is a neighborhood of $x$ and
$\varepsilon>0$, ordered by reverse inclusion in $U$ and by $\geq$ in
$\varepsilon$; this is directed by
\Cref{ex:directed-sets-of-the-book}.  Choose
$y_{U,\varepsilon}\in U$ with
\[
 f(y_{U,\varepsilon})>
 \begin{cases}
 \sup_Uf-\varepsilon, & \sup_Uf<+\infty,\\
 1/\varepsilon, & \sup_Uf=+\infty .
 \end{cases}
\]
The net $y_{U,\varepsilon}$ converges to $x$, because it
lies in $U$ and the neighborhoods are cofinal.  Fix a neighborhood $V$ of $x$
and $\delta>0$.  For $(U,\varepsilon)\geq(V,\delta)$ we have $U\subseteq V$
and $\varepsilon\leq\delta$, hence
\[
 f(y_{U,\varepsilon})>\min\Bigl\{\sup_Uf-\delta,\ 1/\delta\Bigr\}
 \geq\min\Bigl\{f^*(x)-\delta,\ 1/\delta\Bigr\},
\]
using $\sup_Uf\geq f^*(x)$.  Therefore
$\limsup f(y_{U,\varepsilon})\geq\min\{f^*(x)-\delta,1/\delta\}$, and letting
$\delta\downarrow0$ gives $\limsup f(y_{U,\varepsilon})\geq f^*(x)$.  Hence
$L\geq f^*(x)$.

For the final formula, note that $\sup_Uf=\max\{f(x),\sup_{U\setminus\{x\}}f\}$
for every neighborhood $U$ of $x$, and that the infimum over $U$ of a maximum
with the constant $f(x)$ is the maximum of $f(x)$ with the infimum, since
$\Rbar$ is totally ordered and the family is directed.
\end{proof}

If $(f_i)_{i\in I}$ is a family of upper semicontinuous functions, its
pointwise supremum need not be upper semicontinuous.  The natural envelope is
therefore
\[
                       \left(\sup_{i\in I}f_i\right)^*.
\]
This operation occurs in every construction of a psh obstacle envelope.

\begin{example}
\label{ex:usc-spike-regularization}
On $[-1,1]$, let $f_n$ be the continuous triangular function supported in
$[-1/n,1/n]$ with $f_n(0)=1$.  The supremum of the tail family
$\{f_n:n\geq N\}$ is upper semicontinuous, but the pointwise limsup is the
function equal to $1$ at $0$ and $0$ elsewhere.  This function is upper
semicontinuous but not continuous.  Regularization preserves the spike; it
does not mean continuous regularization.
\end{example}

The next two examples are the ones to keep in mind whenever the star in
$(\sup_i\varphi_i)^*$ is about to be dropped.  In the first, the family is an
increasing sequence of continuous functions; in the second, it is an
increasing sequence of subharmonic functions, and the failure locus is exactly
the polar set on which the limit is $-\infty$ before regularization.  The
second is the local model of every phenomenon studied in
\cref{chap:bj-envelope-statement}.

\begin{example}[An increasing family whose supremum is not upper semicontinuous]
\label{ex:sup-not-usc}
On $X=[-1,1]$ put
\[
 f_j(x)=\min\bigl(1,\max(0,jx)\bigr),\qquad j\geq1 .
\]
Each $f_j$ is continuous, hence usc, and $f_j\leq f_{j+1}$: indeed
$\max(0,jx)$ increases in $j$ for $x\geq0$ and vanishes for $x\leq0$.  The
pointwise supremum is
\[
 g=\sup_jf_j=\mathbf 1_{(0,1]} .
\]
This function is not upper semicontinuous: $\{g<1/2\}=[-1,0]$ is not open in
$X$, since $0$ has no neighborhood contained in it.  By
\Cref{prop:usc-limsup-description},
\[
 g^*(0)=\max\Bigl\{g(0),\limsup_{y\to0}g(y)\Bigr\}=\max\{0,1\}=1,
\]
and $g^*=g$ elsewhere, so $g^*=\mathbf 1_{[0,1]}$.  Thus the regularization
changes the function at exactly one point.  Note that
$g^*$ is usc but not continuous, in accordance with the preceding example: the
star is an order-theoretic closure, not a smoothing operator.
\end{example}

\begin{example}[Regularization of an increasing family of subharmonic functions]
\label{ex:sup-not-usc-log}
On the unit disc $D\subset\mathbf C$ put
\[
 u_j(z)=\tfrac1j\log|z|,\qquad j\geq1 .
\]
Each $u_j$ is subharmonic, hence usc, with values in
$\mathbf R\cup\{-\infty\}$, and $u_j\leq u_{j+1}$ because $\log|z|\leq0$ on
$D$.  The pointwise supremum is
\[
 u=\sup_ju_j=
 \begin{cases}
 0,& z\neq0,\\
 -\infty,& z=0,
 \end{cases}
\]
since $\tfrac1j\log|z|\to0$ for $z\neq0$ while $u_j(0)=-\infty$ for every $j$.
This $u$ is not upper semicontinuous at the origin: $\{u<-1\}=\{0\}$ is not
open.  Its regularization is
\[
 u^*\equiv0
\]
by \Cref{prop:usc-limsup-description}, because
$\limsup_{z\to0}u(z)=0$.  The exceptional set $\{u<u^*\}=\{0\}$ is polar,
which is the general phenomenon: for a locally uniformly bounded-above
increasing family of psh functions, the regularization changes the supremum
only on a small set, and quantifying
``small'' is the subject of \cref{sec:capacity-preview} and of
\cref{chap:bfj-capacity}.  Observe also that $u^*=0$ is again subharmonic;
the corresponding statement for arbitrary bounded-above families in the
non-Archimedean setting is precisely the envelope theorem of
\cref{chap:bj-envelope-statement}, and it is not a formality.
\end{example}

The following identities are used constantly and without comment in envelope
computations.  We prove them once.

\begin{lemma}
\label{lem:usc-star-calculus}
Let $X$ be a topological space, let $f,g:X\to\Rbar$, let $(f_i)_{i\in I}$ be a
nonempty family of functions $X\to\Rbar$, let $c\in\mathbf R$, and let
$\lambda>0$.  Then:
\begin{enumerate}
\item $f\leq g$ implies $f^*\leq g^*$, and $(f^*)^*=f^*$;
\item $(f+c)^*=f^*+c$ and $(\lambda f)^*=\lambda f^*$;
\item $\bigl(\sup_if_i\bigr)^*=\bigl(\sup_if_i^*\bigr)^*$;
\item $(\max(f,g))^*=\max(f^*,g^*)$, and more generally the regularization of
  a maximum of finitely many functions is the maximum of the
  regularizations;
\item $(\inf_if_i)^*\leq\inf_if_i^*$, and the inequality can be strict;
\item if $h$ is upper semicontinuous, then $(\sup_i f_i)^*\leq h$ if and only
  if $f_i\leq h$ for every $i$.
\end{enumerate}
\par
\end{lemma}

\begin{proof}
\emph{(1).}  Monotonicity is immediate from \eqref{eq:usc-regularization},
since $\sup_Uf\leq\sup_Ug$ for every $U$.  Idempotence holds because $f^*$ is
usc by \Cref{prop:usc-majorant}, and an usc function equals its own
regularization by the last assertion of that proposition.

\emph{(2).}  For a neighborhood $U$ of $x$ we have
$\sup_U(f+c)=\sup_Uf+c$ and $\sup_U(\lambda f)=\lambda\sup_Uf$ in $\Rbar$,
by \Cref{prop:appE-safe-identities}; taking the infimum over $U$, which
commutes with adding a finite constant and with multiplying by $\lambda>0$
for the same reason, gives the claim.

\emph{(3).}  Since $f_i\leq f_i^*$, part (1) gives ``$\leq$''.  Conversely,
$(\sup_if_i)^*$ is an usc majorant of each $f_i$, hence of each $f_i^*$ by
minimality in \Cref{prop:usc-majorant}, hence of $\sup_if_i^*$; minimality
once more gives $(\sup_if_i^*)^*\leq(\sup_if_i)^*$.

\emph{(4).}  Both sides are usc majorants of $\max(f,g)$: for the right-hand
side, a maximum of two usc functions is usc because
$\{\max(f^*,g^*)<c\}=\{f^*<c\}\cap\{g^*<c\}$ is open.  Minimality gives
$(\max(f,g))^*\leq\max(f^*,g^*)$.  The reverse inequality follows from (1),
since $f\leq\max(f,g)$ gives $f^*\leq(\max(f,g))^*$ and similarly for $g$.
The finite case follows by induction.

\emph{(5).}  From $\inf_if_i\leq f_j$ and (1) we get
$(\inf_if_i)^*\leq f_j^*$ for every $j$.  For strictness take $X=\mathbf R$,
$f_1=\mathbf 1_{\mathbf Q}$ and $f_2=\mathbf 1_{\mathbf R\setminus\mathbf Q}$.
Then $\min(f_1,f_2)=0$, so the left side is $0$, while $f_1^*=f_2^*=1$ and the
right side is $1$.

\emph{(6).}  If $(\sup_if_i)^*\leq h$ then $f_i\leq\sup_if_i\leq h$.
Conversely, if $f_i\leq h$ for all $i$ then $\sup_if_i\leq h$, and $h$ being
usc, minimality in \Cref{prop:usc-majorant} gives $(\sup_if_i)^*\leq h$.
\end{proof}

\begin{remark}
\label{rem:usc-star-not-lattice}
Part (5) is the reason that decreasing regularizations are never written with
a star in this book: for infima, the correct operation is the pointwise
infimum, which is already usc by \Cref{lem:usc-inf}.  Part (6) is the
statement that $f\mapsto(\sup_i f_i)^*$ is the least usc majorant in the
categorical sense, and it is the form in which the envelope property is
verified in \cref{chap:bj-envelope-statement}: one checks the inequality
against every usc competitor rather than computing the star pointwise.
\end{remark}

\begin{lemma}
\label{lem:usc-inf}
The infimum of any family of upper semicontinuous functions is upper
semicontinuous.  In particular, a decreasing net of upper semicontinuous
functions has an upper semicontinuous pointwise limit.
\end{lemma}

\begin{proof}
For $f=\inf_i f_i$ one has
$\{f<c\}=\bigcup_i\{f_i<c\}$, which is open.
\end{proof}

By contrast, the supremum of an \emph{arbitrary} family of usc functions is
usc as soon as the family is finite, by
\Cref{lem:usc-star-calculus}\,(4), and in no other generality: the two
examples above use increasing sequences, which are as well behaved as an
infinite family can be.  The asymmetry between $\sup$ and $\inf$ recorded in
the last two statements is the source of every star in this book.

\subsection{Choquet selection}

The next elementary lemma explains why sequences often reappear after the
topological statements have been made with nets.  It is the single most
frequently used reduction in this book, and it is also the one whose
hypothesis is most often left unchecked in the literature, so we prove it
carefully, isolate the exact property of the space that it uses, and show by
an explicit example that this property cannot be dropped.

\begin{lemma}\label{lem:choquet-selection}
Let $X$ be second countable and let $\mathcal F$ be a nonempty family of
functions $X\to\overline{\mathbf R}$.  There is a countable subfamily
$\{f_j\}_{j\geq1}\subset\mathcal F$ such that
\[
       \left(\sup_{f\in\mathcal F}f\right)^*
       =\left(\sup_{j\geq1}f_j\right)^*.
\]
If $\mathcal F$ is stable under finite maxima, the sequence may be chosen
increasing.
\end{lemma}

\begin{proof}
Choose a countable base $\mathcal B$, and fix $f_0\in\mathcal F$.  For each
$B\in\mathcal B$, set
\[
       M_B=\sup\{f(x):f\in\mathcal F,\ x\in B\}.
\]
If $M_B$ is finite, choose for every $m\geq1$ a function $f_{B,m}$ with
$\sup_Bf_{B,m}>M_B-1/m$; if $M_B=+\infty$, choose
$\sup_Bf_{B,m}>m$.  The cases $M_B=-\infty$ need no choice.  Adjoin $f_0$ to
the selected family.  The resulting family is nonempty and countable and has
the same supremum on every basis element as
$\mathcal F$.  Formula \eqref{eq:usc-regularization} then gives equality of
the regularizations.  If $\mathcal F$ is stable under finite maxima, replace
an enumeration $(g_j)$ of the selected family by
$f_j=g_1\vee\cdots\vee g_j$.

We spell out the two steps that were compressed.  Write
$F=\sup_{f\in\mathcal F}f$ and $G=\sup_jf_j$ for the pointwise suprema of the
full and the selected family, so that $G\leq F$ and hence $G^*\leq F^*$ by
\Cref{lem:usc-star-calculus}\,(1).  For every $B\in\mathcal B$ the two
families have the same supremum over $B$:
\[
 \sup_BG=\sup_j\sup_Bf_j=M_B=\sup_BF,
\]
where the middle equality holds because the selected functions realize $M_B$
in the limit by construction, in the finite case through
$\sup_Bf_{B,m}>M_B-1/m$ and in the infinite case through
$\sup_Bf_{B,m}>m$, while $\sup_Bf_{B,m}\leq M_B$ always.  Since
$\mathcal B$ is a base, every point has a neighborhood basis inside
$\mathcal B$, and \eqref{eq:usc-regularization} may be computed along
$\mathcal B$ alone:
\[
 F^*(x)=\inf_{B\in\mathcal B,\ x\in B}\sup_BF
 =\inf_{B\in\mathcal B,\ x\in B}\sup_BG=G^*(x).
\]
Finally, if $\mathcal F$ is stable under finite maxima then
$f_j=g_1\vee\cdots\vee g_j$ lies in $\mathcal F$, the sequence $(f_j)$ is
increasing, and $\sup_jf_j=\sup_jg_j$, so the conclusion is unchanged.  The
function $f_0$ was adjoined only to guarantee that the selected family is
nonempty when every $M_B$ equals $-\infty$.
\end{proof}

\begin{remark}
\label{rem:choquet-envelope-only}
The lemma does not assert pointwise equality of the two suprema.  It selects
the same upper semicontinuous envelope, which is exactly what is required in
envelope arguments.  Second countability, or a substitute providing a
countable determining family, must be checked rather than assumed.
\end{remark}

\begin{example}[The two suprema really differ]
\label{ex:choquet-pointwise-fails}
Let $X=\mathbf R$, which is second countable, and let $\mathcal F$ be the
family of functions $\mathbf 1_F$ with $F\subset\mathbf R$ finite.  This family
is stable under finite maxima, since
$\mathbf 1_F\vee\mathbf 1_{F'}=\mathbf 1_{F\cup F'}$, and each member is usc.
Its pointwise supremum is the constant $1$.  Choose an increasing sequence
of finite sets $F_j$ whose union is a countable dense subset
$C\subset\mathbf R$, for example $C=\mathbf Q$.  The selected pointwise
supremum is $\mathbf 1_C$, which differs from $1$ at every point of
$\mathbf R\setminus C$.  Nevertheless both regularizations equal the
constant $1$, because every nonempty open set meets this chosen dense set
$C$.  The lemma is therefore sharp: it delivers equality of
regularizations and nothing more.
\end{example}

The proof used second countability only through the following property, and it
is worth isolating it, because several spaces occurring later are not second
countable while still admitting a countable family along which
regularizations may be computed.

\begin{definition}[Countably regularizing family]
\label{def:countably-regularizing}
A countable family $\mathcal B$ of nonempty open subsets of a topological
space $X$ is \emph{regularizing} if for every $f:X\to\Rbar$ and every $x\in X$,
\begin{equation}
\label{eq:countably-regularizing}
 f^*(x)=\inf\bigl\{\textstyle\sup_Bf:\ B\in\mathcal B,\ x\in B\bigr\},
\end{equation}
with the convention $\inf\varnothing=+\infty$ when no member of $\mathcal B$
contains $x$.
\end{definition}

\begin{lemma}\label{lem:regularizing-criterion}
A countable family $\mathcal B$ of nonempty open sets is regularizing if and
only if the members of $\mathcal B$ containing a given point $x$ form a
neighborhood basis at $x$, for every $x\in X$.  In particular a countable base
of the topology is regularizing.
\end{lemma}

\begin{proof}
Suppose the stated basis property holds at $x$.  Every $B\in\mathcal B$
containing $x$ is a neighborhood of $x$, so the right-hand side of
\eqref{eq:countably-regularizing} is at least $f^*(x)$; conversely every
neighborhood $U$ of $x$ contains some $B\in\mathcal B$ with $x\in B$, and then
$\sup_Bf\leq\sup_Uf$, so the right-hand side is at most $f^*(x)$.

Conversely, suppose the basis property fails at some $x$: there is a
neighborhood $U$ of $x$ containing no member of $\mathcal B$ that contains
$x$.  Apply \eqref{eq:countably-regularizing} to
$f=\mathbf 1_{X\setminus U}$.  Then $f(x)=0$, and $f^*(x)=0$ provided $x$ has
a neighborhood, namely the interior of $U$, on which $f$ vanishes; on the
other hand every $B\in\mathcal B$ with $x\in B$ meets $X\setminus U$ by
assumption, so $\sup_Bf=1$ and the right-hand side of
\eqref{eq:countably-regularizing} equals $1$, or equals $+\infty$ if no member
of $\mathcal B$ contains $x$.  In either case it differs from $f^*(x)=0$, so
$\mathcal B$ is not regularizing.
\end{proof}

\begin{proposition}
\label{prop:choquet-regularizing}
Let $X$ be a topological space admitting a countable regularizing family
$\mathcal B$, and let $\mathcal F$ be a nonempty family of functions
$X\to\Rbar$.  Then there is a countable subfamily
$\{f_j\}_{j\geq1}\subseteq\mathcal F$ with
\[
 \Bigl(\sup_{f\in\mathcal F}f\Bigr)^*=\Bigl(\sup_{j\geq1}f_j\Bigr)^* ,
\]
and the sequence may be chosen increasing if $\mathcal F$ is stable under
finite maxima.
\end{proposition}

\begin{proof}
Repeat the proof of \Cref{lem:choquet-selection} verbatim with $\mathcal B$ in
place of a countable base.  Only two properties of $\mathcal B$ were used:
that it is countable, so that the selected family is countable; and that
\eqref{eq:countably-regularizing} computes the regularization, which is the
defining property.  The reduction to an increasing sequence is unchanged.
\end{proof}

The hypothesis is genuinely necessary, and not merely an artifact of the
proof.

\begin{example}[Failure of Choquet selection without countability]
\label{ex:choquet-fails-uncountable}
Let $S$ be an uncountable set and let $X=S\cup\{\infty\}$ be its one-point
compactification, as in \Cref{ex:decreasing-net-not-sequential}: a compact
Hausdorff space.  Let
\[
 \mathcal F=\bigl\{\mathbf 1_F:\ F\subset S\ \text{finite}\bigr\},
\]
a family of usc functions stable under finite maxima, as in
\Cref{ex:choquet-pointwise-fails}.  Its pointwise supremum is
$\mathbf 1_S=\mathbf 1_{X\setminus\{\infty\}}$, whose regularization is the
constant $1$: indeed every neighborhood of $\infty$ is the complement of a
finite subset of $S$ and therefore meets $S$, so the supremum of
$\mathbf 1_S$ over it is $1$.

Let $\{f_j\}_{j\geq1}=\{\mathbf 1_{F_j}\}_{j\geq1}$ be any countable
subfamily and put $C=\bigcup_jF_j$, a countable subset of $S$.  Then
$\sup_jf_j=\mathbf 1_C$.  Choose $s\in S\setminus C$, which exists because $S$
is uncountable.  Since $\{s\}$ is open in $X$ and $\mathbf 1_C(s)=0$,
\[
 \Bigl(\sup_jf_j\Bigr)^*(s)=0\neq1=\Bigl(\sup_{f\in\mathcal F}f\Bigr)^*(s).
\]
So no countable subfamily has the same regularized supremum, and
\Cref{lem:choquet-selection} fails on $X$.  Consistently with
\Cref{lem:regularizing-criterion}, $X$ admits no countable regularizing
family: such a family would have to contain a neighborhood basis at
$\infty$, hence countably many cofinite sets whose complements have countable
union $C$, and then $X\setminus\{s\}$ for $s\in S\setminus C$ would be a
neighborhood of $\infty$ containing no member of the family.
\end{example}

\begin{remark}[Where the countability hypothesis is actually met]
\label{rem:choquet-where-countable}
In the applications the hypothesis is verified in one of three ways.  On a
dual complex, which is a finite polyhedral complex and hence a compact
metrizable space, second countability is automatic; this is the situation of
\cref{chap:skeleta} and of the compactness arguments of
\cref{chap:bfj-compactness}.  On a Berkovich analytification over a field with
a countable dense subfield, the analytification of a variety is metrizable,
and again second countability holds.  In the remaining cases the selection is
performed on a dual complex after the potentials have been shown to be
determined by their restrictions to the skeleta of a cofinal system of
models, using the retractions of \cref{chap:skeleta}; the countable family is
then supplied on the complex, not on the ambient analytification.  What is
never done is to assume metrizability of an arbitrary Berkovich space.
\end{remark}

\section{Radon measures and weak convergence}
\label{sec:radon-weak}

Throughout this section, $X$ is locally compact Hausdorff.  A positive Borel
measure $\mu$ is \emph{Radon} if it is finite on compact sets, inner regular
on open sets, and outer regular on Borel sets.  When $X$ is compact we write
$\mathcal M^+(X)$ for the cone of finite positive Radon measures and
$\mathcal P(X)$ for the probability measures.  The function spaces
$C_c(X)$, $C_b(X)$, and $C(X)$ are real-valued throughout this section.

Spelled out, the two regularity requirements are
\begin{equation}
\label{eq:radon-regularity}
\begin{aligned}
 \mu(U)&=\sup\{\mu(K):K\subseteq U\ \text{compact}\}
 &&(U\ \text{open}),\\
 \mu(B)&=\inf\{\mu(U):U\supseteq B\ \text{open}\}
 &&(B\ \text{Borel}).
\end{aligned}
\end{equation}
On a compact Hausdorff space these conditions are not automatic for a finite
Borel measure, and they are exactly what is used in every approximation
argument below; we therefore never say ``measure'' when we mean ``Radon
measure''.  A useful consequence of \eqref{eq:radon-regularity}, obtained by
complementation inside a compact $X$, is that $\mu$ is also inner regular on
all Borel sets when $\mu(X)<\infty$: given a Borel set $B$ and
$\varepsilon>0$, choose an open $U\supseteq X\setminus B$ with
$\mu(U)\leq\mu(X\setminus B)+\varepsilon$; then $K=X\setminus U$ is a closed,
hence compact, subset of $B$ with $\mu(K)\geq\mu(B)-\varepsilon$.

\subsection{The Riesz representation interface}

Radon measures enter this book only through duality with continuous functions,
and the exact form of that duality is the following imported statement.  We
use it as a black box and state it precisely once.

\begin{theorem}
\label{thm:riesz-interface}
Let $X$ be locally compact Hausdorff and let
$\Lambda:C_c(X,\mathbf R)\to\mathbf R$ be a positive linear functional, that is, a
linear map with $\Lambda(h)\geq0$ whenever $h\geq0$.  Then there is a unique
Radon measure $\mu$ on $X$ with
\begin{equation}
\label{eq:riesz}
 \Lambda(h)=\int_Xh\,\rd\mu\qquad\text{for all }h\in C_c(X,\mathbf R).
\end{equation}
If $X$ is compact, so that $C_c(X,\mathbf R)=C(X,\mathbf R)$, the correspondence
$\mu\mapsto\bigl(h\mapsto\int h\,\rd\mu\bigr)$ is a bijection from
$\mathcal M^+(X)$ onto the positive linear functionals on $C(X,\mathbf R)$, and it
carries $\mathcal P(X)$ onto those with $\Lambda(1)=1$.  Moreover the
continuous linear functionals on $C(X,\mathbf R)$, with the supremum norm, are exactly
the differences of two such integrals, and $\lVert\Lambda\rVert$ is the total
variation of the representing signed measure.
\end{theorem}

\begin{proof}[Foundational measure-theoretic input]
This is \cite[Theorem~2.14]{Rud87} for the existence and uniqueness of $\mu$,
and \cite[Theorem~6.19]{Rud87} for the identification of the dual Banach space
in the compact case; see also \cite[Chapter~7]{Fol99} for the same statements
with the regularity hypotheses displayed in the form
\eqref{eq:radon-regularity}.  We do not reproduce the construction.
\end{proof}

\begin{remark}[What the interface is used for]
\label{rem:riesz-uses}
\Cref{thm:riesz-interface} is invoked in exactly three ways.  First, to
\emph{define} a measure from a functional: this is how Chambert--Loir measures
are constructed in \cref{chap:metrics} and how mixed Monge--Amp\`ere measures
are constructed in \cref{chap:bfj-ma} and \cref{chap:bj-energy}, in each case
by exhibiting a positive linear functional on $C(X)$.  Second, to prove
\emph{uniqueness}: two Radon measures with the same integrals against
continuous functions coincide, which is \Cref{cor:riesz-uniqueness} below.
Third, to obtain \emph{compactness} of $\mathcal P(X)$, in
\Cref{prop:measure-compactness}.  No finer measure theory is required
anywhere in this book, and in particular no completion, no Radon--Nikodym
theorem on a Berkovich space, and no disintegration are used.
\end{remark}

\begin{corollary}\label{cor:riesz-uniqueness}
Let $X$ be compact Hausdorff and let $\mu,\nu\in\mathcal M^+(X)$ satisfy
$\int_Xh\,\rd\mu=\int_Xh\,\rd\nu$ for every $h\in C(X)$.  Then $\mu=\nu$.
\end{corollary}

\begin{proof}
The common functional is positive and linear on $C(X)$, and both $\mu$ and
$\nu$ represent it.  Uniqueness in \Cref{thm:riesz-interface} gives
$\mu=\nu$.  For the reader who prefers a direct argument: by
\eqref{eq:radon-regularity} it suffices to prove $\mu(K)=\nu(K)$ for compact
$K$.  Given $\varepsilon>0$, choose open neighborhoods
$U_\mu,U_\nu\supseteq K$ with
$\mu(U_\mu)\leq\mu(K)+\varepsilon/2$ and
$\nu(U_\nu)\leq\nu(K)+\varepsilon/2$, and put
$U=U_\mu\cap U_\nu$.  Then
$\mu(U)+\nu(U)\leq\mu(K)+\nu(K)+\varepsilon$, and Urysohn's lemma
supplies $h\in C(X)$ with $\mathbf 1_K\leq h\leq\mathbf 1_U$.  Then
\[
 \mu(K)\leq\int h\,\rd\mu=\int h\,\rd\nu\leq\nu(U)\leq\nu(K)+\varepsilon,
\]
and symmetrically.  Let $\varepsilon\downarrow0$.
\end{proof}

For a continuous proper map $p:X\to Y$, the pushforward is
\[
                       p_*\mu(B)=\mu(p^{-1}(B)),
\]
or equivalently
\[
                       \int_Y g\,\rd(p_*\mu)=\int_X g\circ p\,\rd \mu.
\]
Dirac mass at $x$ is denoted by $\delta_x$.

\begin{lemma}\label{lem:pushforward-radon}
Let $p:X\to Y$ be a continuous proper map of locally compact Hausdorff spaces
and let $\mu$ be a Radon measure on $X$.  Then $p_*\mu$ is a Radon measure on
$Y$, the displayed change-of-variables identity holds for every nonnegative
Borel function $g:Y\to[0,+\infty]$ and for every real-valued Borel function
$g$ that is integrable with respect to $p_*\mu$, and
$p_*\delta_x=\delta_{p(x)}$.  If $X$ is compact then
$p_*\mu(Y)=\mu(X)$, so pushforward preserves total mass.
\end{lemma}

\begin{proof}
Borel measurability of $p$ makes $p_*\mu$ a Borel measure, and
$p_*\mu(L)=\mu(p^{-1}L)<\infty$ for compact $L\subseteq Y$ because $p$ is
proper.  For outer regularity, let $B\subseteq Y$ be Borel and
$\varepsilon>0$; outer regularity of $\mu$ gives an open
$V\supseteq p^{-1}(B)$ in $X$ with $\mu(V)\leq\mu(p^{-1}B)+\varepsilon$, but
we need an open set in $Y$.  Properness makes $p$ a closed map, so
$U=Y\setminus p(X\setminus V)$ is open, contains $B$, and satisfies
$p^{-1}(U)\subseteq V$; hence
$p_*\mu(U)\leq\mu(V)\leq p_*\mu(B)+\varepsilon$.  For inner regularity on an
open $U\subseteq Y$, inner regularity of $\mu$ on $p^{-1}(U)$ gives a compact
$K\subseteq p^{-1}(U)$ with $\mu(K)\geq p_*\mu(U)-\varepsilon$, and $p(K)$ is
a compact subset of $U$ with $p_*\mu(p(K))=\mu(p^{-1}p(K))\geq\mu(K)$.

The change-of-variables identity holds for $g=\mathbf 1_B$ by definition,
hence for simple functions by linearity and for nonnegative Borel $g$ by
monotone convergence along an increasing sequence of simple functions.  If
$g$ is integrable with respect to $p_*\mu$, applying the nonnegative identity
to $g^+$ and $g^-$ shows that $g\circ p$ is $\mu$-integrable and gives the
signed identity.  This integrability hypothesis cannot be replaced by
boundedness when $\mu(X)=+\infty$, since both positive and negative parts may
then have infinite integral.  The
identity $p_*\delta_x=\delta_{p(x)}$ is the case $\mu=\delta_x$ of the
definition, and the mass statement is the case $B=Y$.
\end{proof}

\subsection{Weak convergence}

\begin{definition}
\label{def:weak-convergence}
A net $(\mu_a)$ of finite Radon measures on a compact Hausdorff space $X$
\emph{converges weakly} to $\mu$ if
\[
                         \int_X f\,\rd \mu_a\longrightarrow
                         \int_X f\,\rd \mu
\]
for every $f\in C(X)$.
\end{definition}

By \Cref{cor:riesz-uniqueness} the weak limit, if it exists, is unique; the
weak topology on $\mathcal M^+(X)$ is the initial topology for the maps
$\mu\mapsto\int f\,\rd\mu$, $f\in C(X)$, and it is Hausdorff for the same
reason.  Testing against $f=1$ shows that the total mass
$\mu\mapsto\mu(X)$ is weakly continuous on a compact $X$.  On a noncompact
locally compact space the constant $1$ is not a compactly supported test
function, and mass can escape to infinity.  Convergence against $C_c(X)$ is
the \emph{vague} topology.  For finite measures on a noncompact space,
convergence against the bounded continuous functions $C_b(X)$ is the stronger
\emph{narrow}, also commonly
called weak, topology.  We use ``weak'' on compact spaces, ``vague'' for
$C_c$-testing, and ``narrow'' for $C_b$-testing on $\mathbf R$ in
\cref{sec:dh-preview}.

\begin{proposition}\label{prop:measure-compactness}
If $X$ is compact Hausdorff, then $\mathcal P(X)$ is compact for the weak
topology.  If $X$ is compact metrizable, this topology is metrizable and every
sequence in $\mathcal P(X)$ has a weakly convergent subsequence.
\end{proposition}

\begin{proof}
We argue directly in the net language used later.
Consider the map
\[
 \iota:\mathcal P(X)\longrightarrow
 P:=\prod_{f\in C(X)}\bigl[-\lVert f\rVert_\infty,\lVert f\rVert_\infty\bigr],
 \qquad
 \iota(\mu)=\Bigl(\int_Xf\,\rd\mu\Bigr)_{f\in C(X)} .
\]
It is injective by \Cref{cor:riesz-uniqueness}, and by definition of the weak
topology it is a homeomorphism onto its image, the product carrying the
topology of pointwise convergence.  The product $P$ is compact by the
Tychonoff theorem \cite[Chapter~5]{Kel75}.  The image of $\iota$ is closed:
if $\Phi\in P$ lies in the closure, then for $f,g\in C(X)$ and
$s,t\in\mathbf R$ the conditions
\[
 \Phi(sf+tg)=s\Phi(f)+t\Phi(g),
 \qquad
 \Phi(f)\geq0\ \text{ for }f\geq0,
 \qquad
 \Phi(1)=1
\]
each involve only finitely many coordinates and are closed conditions, and
each holds on $\iota(\mathcal P(X))$; hence they hold for $\Phi$.  Thus $\Phi$
is a positive linear functional with $\Phi(1)=1$, and
\Cref{thm:riesz-interface} produces $\mu\in\mathcal P(X)$ with
$\iota(\mu)=\Phi$.  A closed subspace of a compact space is compact, so
$\mathcal P(X)$ is compact, and by \Cref{prop:net-criterion}\,(3) every net of
probability measures has a weakly convergent subnet.

For the metrizable case, let $X$ be compact metrizable.  Then $C(X)$ is
separable in the supremum norm by the Stone--Weierstrass theorem applied to
the algebra generated by a countable family of distance functions
$x\mapsto d(x,x_j)$ for a countable dense set $\{x_j\}$.  Let $(h_k)$ be a
countable dense subset of $C(X)$ and put
\[
 d_{\mathcal P}(\mu,\nu)
 =\sum_{k\geq1}2^{-k}
 \frac{\bigl|\int h_k\,\rd\mu-\int h_k\,\rd\nu\bigr|}
      {1+\lVert h_k\rVert_\infty}.
\]
This is a metric: it vanishes only if $\mu$ and $\nu$ agree against a dense
family of test functions, hence against all of $C(X)$ by uniform
approximation and the uniform mass bound, hence are equal by
\Cref{cor:riesz-uniqueness}.  Its topology is the weak one, again by uniform
approximation.  A compact metric space is sequentially compact, which gives
the last assertion.
\end{proof}

\begin{remark}[Metrizability is a hypothesis, not a convenience]
\label{rem:weak-metrizability}
Only the last sentence of \Cref{prop:measure-compactness} produces
subsequences, and it requires $X$ to be compact metrizable.  For a general
compact Hausdorff $X$ the space $\mathcal P(X)$ is compact but need not be
sequentially compact; the space $X=\{0,1\}^{[0,1]}$ of
\Cref{ex:net-no-subsequence} already gives a sequence $\delta_{x_n}$ with no
weakly convergent subsequence.  Indeed, if
$\delta_{x_{n_j}}\to\mu$, then for $f,g\in C(X)$,
\[
 \int_Xfg\,\rd\mu
 =\lim_j f(x_{n_j})g(x_{n_j})
 =\left(\int_Xf\,\rd\mu\right)
  \left(\int_Xg\,\rd\mu\right).
\]
Thus integration against $\mu$ is a unital multiplicative functional on
$C(X)$.  It follows directly that $\mu$ is a Dirac mass: multiplicativity
with $g=f$ gives zero variance for every $f\in C(X)$, while a Radon
probability measure whose support contained distinct points $x,y$ would admit
an $f\in C(X)$ with $f(x)=0$ and $f(y)=1$, and neighborhoods of $x$ and $y$
of positive $\mu$-mass would make that variance positive.  Hence the support
is a singleton and $\mu=\delta_x$.  Then \Cref{ex:moving-atoms} below would force
$x_{n_j}\to x$, contrary to \Cref{ex:net-no-subsequence}.  This is the reason that
the compactness theorems of \cref{chap:bfj-compactness} and
\cref{chap:bj-strong-topologies} are formulated with nets, and that their
sequential corollaries carry explicit countability hypotheses.
\end{remark}

\begin{proposition}\label{prop:portmanteau}
Let $X$ be compact Hausdorff.  If $\mu_a\to\mu$ weakly in $\mathcal P(X)$, then
\[
 \liminf_a\int_X f\,\rd \mu_a\geq\int_Xf\,\rd \mu
\]
for every lower semicontinuous function $f:X\to\mathbf R\cup\{+\infty\}$
that is bounded below.  Dually,
\[
 \limsup_a\int_X g\,\rd \mu_a\leq\int_Xg\,\rd \mu
\]
for every upper semicontinuous $g:X\to\mathbf R\cup\{-\infty\}$ bounded
above.  Consequently,
\[
 \liminf_a\mu_a(U)\geq\mu(U),\qquad
 \limsup_a\mu_a(F)\leq\mu(F)
\]
for open $U$ and closed $F$.
\end{proposition}

\begin{proof}
A bounded-below lower semicontinuous function is the pointwise supremum of
continuous functions lying below it.  Apply weak convergence to these
continuous minorants and take their supremum.  The upper semicontinuous case
follows by changing signs.  Approximate the indicators of an open set from
below and of a closed set from above.

In detail, let $f$ be lsc with $f\geq m$ for a real constant $m$.  The
integral form of the approximation by continuous minorants,
\[
 \int_Xf\,\rd\nu=\sup\Bigl\{\int_Xh\,\rd\nu:\ h\in C(X),\ h\leq f\Bigr\}
 \qquad(\nu\in\mathcal P(X)),
\]
is proved as \eqref{eq:appE-lsc-integral-minorants}; it uses inner regularity
of $\nu$ and Urysohn's lemma.  For $h\in C(X)$ with $h\leq f$,
\[
 \liminf_a\int_Xf\,\rd\mu_a\geq\liminf_a\int_Xh\,\rd\mu_a
 =\int_Xh\,\rd\mu,
\]
and taking the supremum over $h$ gives the first inequality.  Applying it to
$-g$ gives the second.  For the set statements, let $U$ be open.  The function
$\mathbf 1_U$ is lsc and bounded below, so
$\liminf_a\mu_a(U)\geq\mu(U)$; for closed $F$, the function
$\mathbf 1_F$ is usc and bounded above, so
$\limsup_a\mu_a(F)\leq\mu(F)$.
\end{proof}

\begin{remark}[No convergence on arbitrary Borel sets]
\label{rem:portmanteau-sharp}
Neither inequality in \Cref{prop:portmanteau} can be improved to an equality
for general Borel sets.  On $X=[0,1]$, the measures $\delta_{1/n}$ converge
weakly to $\delta_0$, while
\[
 \delta_{1/n}(\{0\})=0\not\to1=\delta_0(\{0\}),
 \qquad
 \delta_{1/n}\bigl((0,1]\bigr)=1\not\to0=\delta_0\bigl((0,1]\bigr).
\]
The first displays failure for a closed set on which the limsup inequality is
strict, and the second failure for an open set on which the liminf inequality
is strict.  Equality does hold for a Borel set $B$ whose topological boundary
is $\mu$-null, by applying both inequalities to $\Int B$ and $\overline B$;
this is the only form of set convergence used in this book, and it is invoked
only after the boundary has been shown to be negligible.
\end{remark}

\begin{example}[Moving atoms]
\label{ex:moving-atoms}
Let $X$ be compact Hausdorff.  If $x_a\to x$ in $X$, then
$\delta_{x_a}\to\delta_x$ weakly.  Conversely, if
$\delta_{x_a}\to\delta_x$, then $x_a\to x$:
continuous functions separate $x$ from every closed set not containing it.
Thus the Dirac embedding $x\mapsto\delta_x$ is a topological embedding.
\end{example}

For the converse direction the argument is the following.  Suppose
$x_a\not\to x$; then some open $U\ni x$ fails to contain a tail, so the net is
frequently in the closed set $F=X\setminus U$.  A compact Hausdorff space is
normal, so Urysohn's lemma gives $h\in C(X)$ with $h(x)=1$ and $h\equiv0$ on
$F$.  Then $\int h\,\rd\delta_{x_a}=h(x_a)$ is frequently $0$, while
$\int h\,\rd\delta_x=1$; hence $\delta_{x_a}\not\to\delta_x$.  Injectivity of
$x\mapsto\delta_x$ follows from the same separation property, and continuity
in the forward direction is \Cref{prop:net-criterion}\,(2) applied to each
$h\in C(X)$.

\begin{example}[Escape of mass in the noncompact case]
\label{ex:vague-escape}
On $X=\mathbf R$, which is locally compact Hausdorff but not compact, the
measures $\delta_n$ satisfy $\int h\,\rd\delta_n=h(n)\to0$ for every
$h\in C_c(\mathbf R)$, so $\delta_n\to0$ vaguely, while $\delta_n(\mathbf R)=1$
for every $n$.  Total mass is therefore not vaguely continuous, and a vague
limit of probability measures need not be a probability measure.  Compactifying
$\mathbf R$ to $[-\infty,+\infty]$ repairs this: there $\delta_n\to
\delta_{+\infty}$ weakly, and no mass is lost.  The corresponding phenomenon
in this book is the compactness of $X^{\An}$ for proper $X$, established in
\cref{chap:berkovich}; it is what allows every Monge--Amp\`ere measure to be
treated as a measure of fixed total mass.
\end{example}

The following two statements are the ones actually applied when a
Monge--Amp\`ere measure and a potential vary simultaneously.

\begin{lemma}
\label{lem:joint-weak-uniform}
Let $X$ be compact Hausdorff, let $\mu_a\to\mu$ weakly in
$\mathcal M^+(X)$ with $\sup_a\mu_a(X)=:C<\infty$, and let $(g_a)$ be a net in
$C(X)$, indexed by the same directed set, converging uniformly to $g\in C(X)$.
Then
\[
 \int_Xg_a\,\rd\mu_a\longrightarrow\int_Xg\,\rd\mu .
\]
\end{lemma}

\begin{proof}
Write
\[
 \int_Xg_a\,\rd\mu_a-\int_Xg\,\rd\mu
 =\int_X(g_a-g)\,\rd\mu_a+\Bigl(\int_Xg\,\rd\mu_a-\int_Xg\,\rd\mu\Bigr).
\]
The first term is bounded in absolute value by
$C\lVert g_a-g\rVert_\infty$, which tends to $0$; the second tends to $0$ by
weak convergence.  Given $\varepsilon>0$, choose an index beyond which each
term is at most $\varepsilon/2$ in absolute value, using directedness to
combine the two thresholds.
\end{proof}

\begin{lemma}
\label{lem:usc-along-decreasing-net}
Let $X$ be compact Hausdorff, let $\mu_a\to\mu$ weakly in $\mathcal P(X)$, and
let $(g_b)_{b\in B}$ be a decreasing net of upper semicontinuous functions
$X\to\mathbf R\cup\{-\infty\}$, uniformly bounded above, with
$g=\inf_bg_b$.  Then
\[
 \limsup_a\int_Xg\,\rd\mu_a\leq\int_Xg\,\rd\mu ,
\]
and moreover
\[
 \inf_b\ \limsup_a\int_Xg_b\,\rd\mu_a\leq\int_Xg\,\rd\mu .
\]
\end{lemma}

\begin{proof}
The first inequality is \Cref{prop:portmanteau} applied to the usc function
$g$, which is usc by \Cref{lem:usc-inf} and bounded above by hypothesis.  For
the second, \Cref{prop:portmanteau} applied to each $g_b$ gives
$\limsup_a\int g_b\,\rd\mu_a\leq\int g_b\,\rd\mu$; taking the infimum over $b$
and using \Cref{cor:usc-net-monotone-integral} for the measure $\mu$ yields
\[
 \inf_b\limsup_a\int_Xg_b\,\rd\mu_a
 \leq\inf_b\int_Xg_b\,\rd\mu=\int_Xg\,\rd\mu .
\]
\end{proof}

\begin{remark}
\label{rem:no-reverse-inequality}
The inequalities of \Cref{prop:portmanteau} and
\Cref{lem:usc-along-decreasing-net} are strict in general, and no reverse
inequality is available.  On $X=[0,1]$ take $\mu_m=\delta_{1/m}$, which
converges weakly to $\mu=\delta_0$, and $g=\mathbf 1_{\{0\}}$, which is upper
semicontinuous and bounded.  Then
\[
 \limsup_m\int_Xg\,\rd\mu_m=0<1=\int_Xg\,\rd\mu .
\]
Every place in this book where a \emph{limit} of integrals is asserted,
rather than an inequality, therefore rests on a separate continuity theorem:
uniform convergence, as in \Cref{lem:joint-weak-uniform}; monotone convergence
with a fixed measure, as in \Cref{cor:usc-net-monotone-integral}; or a
capacity estimate, as in \cref{chap:bfj-capacity}.
\end{remark}

\subsection{Concave energy and its dual}

The variational theories later in the book fit the following elementary
pattern.  Let $\mathcal H$ be a convex subset of $C(X)$ for a compact
Hausdorff space $X$, stable under addition of constants, and let
$E:\mathcal H\to\mathbf R\cup\{-\infty\}$ be concave.  Assume that for a
fixed number $M>0$,
\begin{equation}
\label{eq:energy-translation}
                         E(\varphi+c)=E(\varphi)+Mc.
\end{equation}
For a finite positive measure $\mu$, define the dual energy
\begin{equation}
\label{eq:dual-energy-abstract}
 E^\vee(\mu)=
 \sup_{\varphi\in\mathcal H}
 \left(E(\varphi)-\int_X\varphi\,\rd \mu\right).
\end{equation}

\begin{lemma}\label{lem:dual-mass}
If \eqref{eq:energy-translation} holds and $E(\varphi_0)$ is finite for some
$\varphi_0\in\mathcal H$, then $E^\vee(\mu)<+\infty$ is possible only when
$\mu(X)=M$.
\end{lemma}

\begin{proof}
Fix such a $\varphi_0$ and replace it by $\varphi_0+c$.  The expression in
\eqref{eq:dual-energy-abstract} changes by
$c(M-\mu(X))$.  Its supremum over $c\in\mathbf R$ is infinite unless the
coefficient vanishes.
\end{proof}

For a measure of the correct mass, the definition immediately gives the
Fenchel inequality
\begin{equation}
\label{eq:fenchel-energy}
                E(\varphi)-\int_X\varphi\,\rd \mu\leq E^\vee(\mu).
\end{equation}
Equality is the variational form of the Monge--Amp\`ere equation.  No
biconjugacy statement is implicit here; such a statement requires the
topology, semicontinuity, and compactness hypotheses proved later in each
particular theory.

\subsection{Legendre duality in finite dimension}
\label{subsec:legendre-finite}

The reason for the caution in the last paragraph becomes visible once the
finite-dimensional model is written out.  There the duality is complete, and
we prove it, because every later variational statement is an attempt to
reproduce it in a space of potentials.

Let $V$ be a finite-dimensional real vector space with dual $V^\vee$ and
pairing $\langle\cdot,\cdot\rangle$.  For
$F:V\to\mathbf R\cup\{-\infty\}$ concave, define its \emph{concave conjugate}
\begin{equation}
\label{eq:concave-conjugate}
 F^\vee(\ell)=\sup_{v\in V}\bigl(F(v)-\langle\ell,v\rangle\bigr),
 \qquad \ell\in V^\vee ,
\end{equation}
a function $V^\vee\to\mathbf R\cup\{+\infty\}$.  This is the sign convention
of \eqref{eq:dual-energy-abstract}: $F$ plays the role of the energy $E$, the
variable $v$ the role of the potential $\varphi$, and $\ell$ the role of the
measure $\mu$ through $\varphi\mapsto\int\varphi\,\rd\mu$.

\begin{proposition}
\label{prop:concave-biconjugation}
Let $F:V\to\mathbf R\cup\{-\infty\}$ be concave and not identically
$-\infty$.  Then:
\begin{enumerate}
\item $F^\vee$ is convex and lower semicontinuous, and
\begin{equation}
\label{eq:fenchel-young-concave}
 F(v)-\langle\ell,v\rangle\leq F^\vee(\ell)
 \qquad(v\in V,\ \ell\in V^\vee);
\end{equation}
\item $F^{\vee\vee}(v):=\inf_{\ell\in V^\vee}
  \bigl(F^\vee(\ell)+\langle\ell,v\rangle\bigr)$
  is the smallest upper semicontinuous concave majorant of $F$; in particular
  $F^{\vee\vee}=F$ if and only if $F$ is upper semicontinuous;
\item equality holds in \eqref{eq:fenchel-young-concave} for a given pair
  $(v,\ell)$ if and only if $\ell$ is a supergradient of $F$ at $v$, that is
\[
 F(w)\leq F(v)+\langle\ell,w-v\rangle\qquad\text{for all }w\in V .
\]
\end{enumerate}
\par
\end{proposition}

\begin{proof}
Put $f=-F$, a convex function $V\to\mathbf R\cup\{+\infty\}$ that is proper in
the sense of \Cref{def:appD-proper-closed}, since $F\not\equiv-\infty$ and $F$
never takes the value $+\infty$.  Then
\[
 F^\vee(\ell)=\sup_v\bigl(-f(v)-\langle\ell,v\rangle\bigr)
 =\sup_v\bigl(\langle-\ell,v\rangle-f(v)\bigr)=f^*(-\ell),
\]
with $f^*$ the convex conjugate \eqref{eq:appD-conjugate}.  Assertion (1) is
therefore the statement that $f^*$ is convex and lsc, which holds because
$f^*$ is a supremum of continuous affine functions, together with the
Fenchel--Young inequality \eqref{eq:appD-fenchel} of
\Cref{prop:appD-fenchel-equality}.

For (2), the same substitution gives
\[
 F^{\vee\vee}(v)=\inf_\ell\bigl(f^*(-\ell)+\langle\ell,v\rangle\bigr)
 =-\sup_{m\in V^\vee}\bigl(\langle m,v\rangle-f^*(m)\bigr)
 =-f^{**}(v),
\]
after the substitution $m=-\ell$.  The Fenchel--Moreau theorem
\Cref{thm:appD-fenchel-moreau} identifies $f^{**}$ with the lower
semicontinuous convex closure $\overline f$, that is, the largest lsc convex
minorant of $f$.  Negating, $F^{\vee\vee}$ is the smallest usc concave
majorant of $F$, and it equals $F$ exactly when $F$ is usc.

For (3), equality in \eqref{eq:fenchel-young-concave} says
$F(v)-\langle\ell,v\rangle=\sup_w(F(w)-\langle\ell,w\rangle)$, that is,
$F(w)-\langle\ell,w\rangle\leq F(v)-\langle\ell,v\rangle$ for every $w$, which
rearranges to the supergradient inequality.
\end{proof}

\begin{remark}[What the infinite-dimensional theory must supply]
\label{rem:legendre-infinite-dimensional}
\Cref{prop:concave-biconjugation} uses finite-dimensionality only through
\Cref{thm:appD-fenchel-moreau}, whose proof separates a point from a closed
convex epigraph.  In the variational chapters the role of $V$ is played by a
space of potentials, which is not a finite-dimensional vector space and is not
even a topological vector space in any useful sense; the separation step must
then be replaced by a compactness and semicontinuity argument specific to the
theory at hand.  This is why \eqref{eq:fenchel-energy} is stated as an
inequality and why the corresponding equality is proved only after
differentiability of the energy, in \cref{chap:bfj-differentiability} and
\cref{chap:bj-finite-energy}.  The finite-dimensional statement is the
template, not a lemma that can be quoted.
\end{remark}

The next example is a complete computation of a Legendre pair.  It is the
finite model of the free-energy/entropy duality that reappears as the
Ding and Mabuchi functionals in \cref{chap:bj-stability}, and it exhibits both
the translation identity \eqref{eq:energy-translation} and the mass constraint
of \Cref{lem:dual-mass}.

\begin{example}[A Legendre pair: free energy and relative entropy]
\label{ex:legendre-entropy-pair}
Let $X=\{1,\dots,N\}$ be a finite set, let $p=(p_1,\dots,p_N)$ be a
probability measure on $X$ with $p_i\geq0$ and $\sum_ip_i=1$, and let
$\mathcal H=\mathbf R^N$ be the space of all real functions
$\varphi=(\varphi_i)$ on $X$.  Define the \emph{free energy}
\[
 E(\varphi)=-\log Z(\varphi),
 \qquad
 Z(\varphi)=\sum_{i=1}^Np_i\rme^{-\varphi_i}.
\]
We compute $E^\vee$ completely.

\emph{Step 1: $E$ is concave and satisfies \eqref{eq:energy-translation} with
$M=1$.}  The function $\psi\mapsto\log\sum_ip_i\rme^{\psi_i}$ is convex, by
H\"older's inequality: for $s\in[0,1]$,
\[
 \sum_ip_i\rme^{(1-s)\psi_i+s\chi_i}
 \leq\Bigl(\sum_ip_i\rme^{\psi_i}\Bigr)^{1-s}
     \Bigl(\sum_ip_i\rme^{\chi_i}\Bigr)^{s}.
\]
Hence $E(\varphi)=-\log\sum_ip_i\rme^{-\varphi_i}$ is concave in $\varphi$.
Replacing $\varphi$ by $\varphi+c$ multiplies $Z$ by $\rme^{-c}$, so
$E(\varphi+c)=E(\varphi)+c$.

\emph{Step 2: the mass constraint.}  By \Cref{lem:dual-mass}, applied with
$M=1$ and $\varphi_0=0$, for which $E(0)=0$ is finite, the dual energy
$E^\vee(\mu)$ is $+\infty$ unless $\mu(X)=\sum_i\mu_i=1$.  We may therefore
assume that $\mu$ is a probability measure on $X$.

\emph{Step 3: the key identity.}  Fix $\varphi\in\mathbf R^N$ and set
\[
 \nu_i=\frac{p_i\rme^{-\varphi_i}}{Z(\varphi)},
\]
the associated Gibbs measure, a probability measure on $X$ with $\nu_i=0$
exactly when $p_i=0$.  Let $\mu$ be a probability measure with
$\mu_i=0$ whenever $p_i=0$, and write
\[
 H(\mu\mid\lambda)=\sum_{i:\,\mu_i>0}\mu_i\log\frac{\mu_i}{\lambda_i}
 \in[0,+\infty]
\]
for the relative entropy of $\mu$ with respect to a probability measure
$\lambda$, with the convention that the sum is $+\infty$ if $\lambda_i=0$ for
some $i$ with $\mu_i>0$.  For $i$ with $\mu_i>0$ we have $p_i>0$ and
$\nu_i>0$, and
$\varphi_i=\log p_i-\log\nu_i-\log Z(\varphi)$; therefore
\[
 E(\varphi)-\sum_i\varphi_i\mu_i
 =-\log Z-\sum_{i:\,\mu_i>0}\mu_i
   \bigl(\log p_i-\log\nu_i-\log Z\bigr)
 =\sum_{i:\,\mu_i>0}\mu_i\log\frac{\nu_i}{p_i},
\]
using $\sum_i\mu_i=1$.  Adding and subtracting $\mu_i\log\mu_i$ gives
\begin{equation}
\label{eq:gibbs-identity}
 E(\varphi)-\sum_i\varphi_i\mu_i
 =H(\mu\mid p)-H(\mu\mid\nu).
\end{equation}

\emph{Step 4: nonnegativity of relative entropy.}  For probability measures
$\mu,\lambda$ with $\mu\ll\lambda$, Jensen's inequality applied to the strictly
convex function $t\mapsto t\log t$ and the measure $\lambda$ gives
\[
 H(\mu\mid\lambda)=\sum_i\lambda_i\,
 \frac{\mu_i}{\lambda_i}\log\frac{\mu_i}{\lambda_i}
 \geq\Bigl(\sum_i\mu_i\Bigr)\log\Bigl(\sum_i\mu_i\Bigr)=0,
\]
with equality if and only if $\mu_i/\lambda_i$ is constant, that is
$\mu=\lambda$.

\emph{Step 5: conclusion.}  Combining \eqref{eq:gibbs-identity} with Step 4,
\[
 E(\varphi)-\sum_i\varphi_i\mu_i\leq H(\mu\mid p)
\]
for every $\varphi$, with equality if and only if $\mu=\nu$, that is, if and
only if $\varphi_i=-\log(\mu_i/p_i)+\mathrm{const}$ on $\{p_i>0\}$.  Such a
$\varphi$ exists precisely when $\mu_i>0$ for every $i$ with $p_i>0$.  In
general, take $\varphi_i=-\log(\mu_i/p_i)$ where $\mu_i>0$ and
$\varphi_i=T>0$ elsewhere; then $Z=1+\rme^{-T}\sum_{i:\,\mu_i=0}p_i$, and
$\nu_i=\mu_i/Z$ for $\mu_i>0$, so $H(\mu\mid\nu)=\log Z\to0$ as
$T\to+\infty$.  Hence the supremum equals $H(\mu\mid p)$ in all cases.
Finally, if $\mu_i>0=p_i$ for some $i$, let $\varphi_i\to-\infty$ at that
index alone: $Z$ is unchanged, while $-\sum_j\varphi_j\mu_j\to+\infty$, so
$E^\vee(\mu)=+\infty=H(\mu\mid p)$.  We conclude
\begin{equation}
\label{eq:free-energy-dual}
 E^\vee(\mu)=
 \begin{cases}
 H(\mu\mid p), & \mu\ \text{a probability measure on }X,\\
 +\infty, & \text{otherwise}.
 \end{cases}
\end{equation}

\emph{Step 6: biconjugation.}  Since $E$ is concave and continuous on
$\mathbf R^N$, \Cref{prop:concave-biconjugation}\,(2) gives
$E^{\vee\vee}=E$, that is
\[
 -\log\sum_ip_i\rme^{-\varphi_i}
 =\inf_{\mu\in\mathcal P(X)}
 \Bigl(H(\mu\mid p)+\sum_i\varphi_i\mu_i\Bigr),
\]
which is the Gibbs variational principle.  The minimizer is the Gibbs measure
$\nu$ of Step 3, in agreement with
\Cref{prop:concave-biconjugation}\,(3): the function $E$ is differentiable
with
\[
 \frac{\partial E}{\partial\varphi_i}(\varphi)
 =-\frac1{Z(\varphi)}\frac{\partial Z}{\partial\varphi_i}(\varphi)
 =\frac{p_i\rme^{-\varphi_i}}{Z(\varphi)}=\nu_i ,
\]
so its unique supergradient at $\varphi$ is the linear form
$\psi\mapsto\sum_i\psi_i\nu_i$, that is, the measure $\nu$.  Consistently with
Step 1, the derivatives sum to $1=M$.

The dictionary with the later chapters is exact: $E$ is the energy, $E^\vee$
is the entropy of a measure relative to a reference, the mass constraint of
Step 2 is the statement that only measures of total mass $(L^n)$ have finite
dual energy, and the equality case of Step 5 is the Euler--Lagrange equation
of the variational problem, that is, the Monge--Amp\`ere equation.
\end{example}

\begin{example}[A one-variable Legendre pair and its inverse]
\label{ex:legendre-one-variable}
The transform used for test curves in \cref{chap:rays-test-curves} is
\eqref{eq:appD-concave-transform}:
$\mathcal Lg(t)=\sup_{\tau\in\mathbf R}\{g(\tau)+t\tau\}$ for $t\geq0$.  Take
\[
 g(\tau)=
 \begin{cases}
 0,&\tau\leq0,\\
 -\tau,&0\leq\tau\leq1,\\
 -\infty,&\tau>1 .
 \end{cases}
\]
This $g$ is concave, upper semicontinuous, and nonincreasing, so
\Cref{thm:appD-concave-duality} applies.  We compute both directions
explicitly.

For $t\geq0$, split the supremum.  On $\tau\leq0$ the expression $g+t\tau$
equals $t\tau\leq0$, with supremum $0$ attained at $\tau=0$.  On
$0\leq\tau\leq1$ it equals $(t-1)\tau$, with supremum $\max\{0,t-1\}$.  On
$\tau>1$ it is $-\infty$.  Hence
\[
 \mathcal Lg(t)=\max\{0,t-1\}=(t-1)_+ ,
\]
a convex, piecewise affine, nondecreasing function with a single break at
$t=1$ and slopes $0$ and $1$; these are exactly the endpoints of the interval
$[0,1]$ on which $g$ is finite and affine.

For the inverse transform, compute
$\inf_{t\geq0}\{(t-1)_+-t\tau\}$.  If $\tau\leq0$, both terms are
nonnegative, and the value $0$ is attained at $t=0$; the infimum is
$0=g(\tau)$.  If $0\leq\tau\leq1$, then on $0\leq t\leq1$ the expression is
$-t\tau$, decreasing, with minimum $-\tau$ at $t=1$, and on $t\geq1$ it is
$t(1-\tau)-1$, nondecreasing, with value $-\tau$ at $t=1$; the infimum is
$-\tau=g(\tau)$.  If $\tau>1$, then on $t\geq1$ the expression
$t(1-\tau)-1$ tends to $-\infty$; the infimum is $-\infty=g(\tau)$.  Thus
\eqref{eq:appD-concave-inversion} holds, as it must.

Two features of this computation recur.  First, the break points of
$\mathcal Lg$ record the slopes of $g$, and conversely; this is the mechanism
by which a test curve determines a ray and vice versa.  Second, the value
$-\infty$ of $g$ for $\tau>1$ is what makes $\mathcal Lg$ finite, and dropping
it would make $\mathcal Lg\equiv+\infty$ for $t>1$.  Extended-real values are
not a blemish on the theory; they carry the information.
\end{example}

\section{Capacities, quasicontinuity, and polar sets: a preview}
\label{sec:capacity-preview}

Measures are well suited to integration, but sets negligible for one measure
need not be negligible for all Monge--Amp\`ere measures.  A capacity records
uniform negligibility.

This section is axiomatic on purpose.  In \cref{chap:bfj-capacity} and
\cref{chap:bj-strong-topologies} several different capacities are constructed
from Monge--Amp\`ere operators, and each of them has its own analytic
input; what the later chapters import from here is only the formal calculus of
a capacity, its polar sets, and its quasicontinuous functions.  Isolating that
calculus once prevents the same argument from being rewritten in each theory,
and it makes visible which of the axioms each application actually uses.  In
particular, countable subadditivity is stated separately every time, because
it is the axiom that fails most easily and the one on which the useful
statements depend.

\begin{definition}
\label{def:choquet-capacity}
Let $X$ be a topological space.  A \emph{Choquet capacity} is a function
\[
                 \operatorname{Cap}:2^X\longrightarrow[0,+\infty]
\]
such that:
\begin{enumerate}
\item $\operatorname{Cap}(\varnothing)=0$ and
  $A\subset B$ implies $\operatorname{Cap}(A)\leq\operatorname{Cap}(B)$;
\item if $A_j\nearrow A$, then
  $\operatorname{Cap}(A_j)\nearrow\operatorname{Cap}(A)$;
\item if $K_j\searrow K$ are compact, then
  $\operatorname{Cap}(K_j)\searrow\operatorname{Cap}(K)$.
\end{enumerate}
We will additionally state countable subadditivity whenever it is available;
it is not part of every author's minimal definition of a Choquet capacity.
\end{definition}

A set $P$ is \emph{polar} for the capacity if
$\operatorname{Cap}(P)=0$.  Equip $\Rbar$ with the metric
\[
 \rho(s,t)=|\vartheta(s)-\vartheta(t)|,
 \qquad
 \vartheta(t)=\frac{2}{\pi}\arctan t,
 \qquad
 \vartheta(\pm\infty)=\pm1.
\]
This metric induces the order topology.  Throughout this section, uniform
convergence of $\Rbar$-valued functions is taken with respect to $\rho$; on
functions with values in a fixed compact interval of $\mathbf R$, it is
equivalent to ordinary uniform convergence.  A function
$f:X\to\overline{\mathbf R}$ is
\emph{quasicontinuous} if for every $\varepsilon>0$ there is an open set $U$
with $\operatorname{Cap}(U)<\varepsilon$ such that
$f|_{X\setminus U}$ is continuous.

Three remarks on the definition.  The continuity required of
$f|_{X\setminus U}$ is continuity of the restriction for the subspace
topology, which is weaker than continuity of $f$ at the points of
$X\setminus U$.  The exceptional set is required to be \emph{open}; requiring
only a Borel set of small capacity gives a formally weaker notion, and the two
agree whenever $\Capa$ is outer regular in the sense that
$\Capa(A)=\inf\{\Capa(U):U\supseteq A\ \text{open}\}$, which holds for every
capacity constructed later.  Finally, a quasicontinuous function may take the
values $\pm\infty$; on the complement of the exceptional set it is then
continuous as a map into $\Rbar$ with its order topology.

\begin{example}[A Choquet capacity that is not countably subadditive]
\label{ex:capacity-not-subadditive}
Let $X=\{0,1\}$ with the discrete topology and define
$\Capa(A)=|A|^2$.  Monotonicity and normalization are immediate.  Every
increasing or decreasing sequence of subsets is eventually constant, so
axioms (2) and (3) hold.  Thus $\Capa$ is a Choquet capacity.  It is not
subadditive, since
\[
 \Capa(\{0,1\})=4>2=\Capa(\{0\})+\Capa(\{1\}).
\]
The point is that countable subadditivity does not follow from axioms
(1)--(3).  Whenever a capacity is
constructed in this book, countable subadditivity is proved, not assumed.
\end{example}

\begin{proposition}
\label{prop:polar-sigma-ideal}
Let $\Capa$ be a countably subadditive Choquet capacity on $X$.  Then:
\begin{enumerate}
\item a subset of a polar set is polar;
\item a countable union of polar sets is polar;
\item if $\mu$ is a positive Borel measure on $X$ which is inner regular on
  Borel sets and satisfies $\mu(K)\leq C\,\Capa(K)$ for every compact
  $K\subseteq X$ and some constant $C$, then every Borel polar set is
  $\mu$-null.
\end{enumerate}
\par
\end{proposition}

\begin{proof}
\emph{(1)} is monotonicity, axiom (1).  \emph{(2)} If
$\Capa(P_j)=0$ for $j\geq1$, then
$\Capa\bigl(\bigcup_jP_j\bigr)\leq\sum_j\Capa(P_j)=0$.
\emph{(3)} Let $P$ be Borel with $\Capa(P)=0$.  For compact $K\subseteq P$,
monotonicity gives $\Capa(K)=0$, hence $\mu(K)=0$.  Inner regularity gives
$\mu(P)=\sup\{\mu(K):K\subseteq P\ \text{compact}\}=0$.
\end{proof}

\begin{remark}[The interface used later]
\label{rem:capacity-interface}
Part (3) of \Cref{prop:polar-sigma-ideal} is the exact form in which
capacities are used to compare null sets.  In \cref{chap:bfj-capacity} the
measures $\mu$ are Monge--Amp\`ere measures of bounded potentials, the
constant $C$ comes from a Chern--Levine--Nirenberg type estimate, and the
conclusion is that a bounded potential is determined outside a polar set.  The
hypothesis of inner regularity is automatic for Radon measures by
\eqref{eq:radon-regularity}.  No statement about \emph{non}polar sets is
imported from this section; those require the specific capacity.
\end{remark}

\begin{proposition}
\label{prop:quasicontinuity-calculus}
Let $\Capa$ be a countably subadditive Choquet capacity on $X$.
\begin{enumerate}
\item If $f_1,\dots,f_k$ are quasicontinuous, then for every $\varepsilon>0$
  there is a single open $U$ with $\Capa(U)<\varepsilon$ such that every
  $f_i|_{X\setminus U}$ is continuous.  The same holds for a countable family.
\item Finite sums of $\mathbf R\cup\{-\infty\}$-valued
  quasicontinuous functions, finite maxima, finite minima, and products of
  real-valued quasicontinuous functions are quasicontinuous, as is
  $\chi\circ f$ for $\chi:\Rbar\to\Rbar$ continuous.
\item A $\rho$-uniform limit of quasicontinuous functions is
  quasicontinuous.
\item If $f_j$ is quasicontinuous for each $j$ and $f_j\to f$ uniformly
  outside a set of arbitrarily small capacity, in the sense that for every
  $\varepsilon>0$ there is an open $U$ with $\Capa(U)<\varepsilon$ and
  $f_j\to f$ $\rho$-uniformly on $X\setminus U$, then $f$ is
  quasicontinuous.
\end{enumerate}
\par
\end{proposition}

\begin{proof}
\emph{(1).}  Given $\varepsilon>0$ and a countable family $(f_j)_{j\geq1}$,
choose open $U_j$ with $\Capa(U_j)<2^{-j}\varepsilon$ such that
$f_j|_{X\setminus U_j}$ is continuous, and put $U=\bigcup_jU_j$, which is open
with $\Capa(U)\leq\sum_j\Capa(U_j)<\varepsilon$ by countable subadditivity.
For each $j$, the set $X\setminus U$ is a subspace of $X\setminus U_j$, and a
restriction of a continuous map to a subspace is continuous.

\emph{(2).}  By (1) there is one open $U$ of capacity $<\varepsilon$ outside
which all the given functions are continuous.  Addition on
$\mathbf R\cup\{-\infty\}$, with $-\infty$ absorbing, and maximum and minimum
on $\Rbar$ are continuous for the order topology.  Products of continuous
real-valued functions on the subspace $X\setminus U$ are continuous there,
and so is $\chi\circ f$ for continuous $\chi$.

\emph{(3).}  Let $f_j\to f$ uniformly with each $f_j$ quasicontinuous.  By (1)
choose an open $U$ with $\Capa(U)<\varepsilon$ outside which every $f_j$ is
continuous.  On $X\setminus U$, the function $f$ is a uniform limit of
continuous functions, hence continuous.

\emph{(4).}  Given $\varepsilon>0$, choose an open $U'$ with
$\Capa(U')<\varepsilon/2$ on whose complement $f_j\to f$ uniformly, and by (1)
an open $U''$ with $\Capa(U'')<\varepsilon/2$ on whose complement every $f_j$
is continuous.  Then $U=U'\cup U''$ is open with $\Capa(U)<\varepsilon$, and
on $X\setminus U$ the function $f$ is a uniform limit of continuous functions,
hence continuous.
\end{proof}

\begin{remark}[What is not proved here]
\label{rem:quasicontinuity-limits}
Two statements that look similar to
\Cref{prop:quasicontinuity-calculus} are \emph{not} available at this level of
generality, and each is proved separately in the theory where it is used.
The first is that a decreasing limit of quasicontinuous functions is
quasicontinuous; the proof requires convergence in capacity, not merely
pointwise convergence.  The second is that two quasicontinuous functions
agreeing outside a polar set agree quasi-everywhere in the strong sense that
their difference has quasi-open sublevel sets; this needs a quasi-Lindel\"of
property of the capacity, which holds for the Monge--Amp\`ere capacities of
\cref{chap:bfj-capacity} but is not a consequence of axioms (1)--(3) together
with countable subadditivity.  We do not use either statement before it has
been proved for the capacity at hand.
\end{remark}

\begin{definition}
\label{def:capacity-convergence}
Let the capacity be finite, meaning $\operatorname{Cap}(X)<+\infty$, let
$(Y,d)$ be a metric space, and let
$f_j,f:X\to Y$.  The sequence $f_j$ \emph{converges to $f$ in capacity} if
for every $\varepsilon>0$,
\[
 \operatorname{Cap}\bigl(\{d(f_j,f)>\varepsilon\}\bigr)\longrightarrow0.
\]
For real-valued functions we take $d(s,t)=|s-t|$.  For $\Rbar$-valued
functions we may take the metric $\rho$ above, so no undefined difference of
infinite values occurs.  On uniformly bounded real-valued families the two
choices give the same notion of convergence.
\end{definition}

\begin{proposition}\label{prop:capacity-subsequence}
Assume that $\operatorname{Cap}$ is countably subadditive.  If
$f_j:X\to(Y,d)$ converges in capacity to $f:X\to(Y,d)$, then some
subsequence converges pointwise to $f$ outside a polar set.
\end{proposition}

\begin{proof}
Choose $j_m$ so that
\[
 \operatorname{Cap}\bigl(\{d(f_{j_m},f)>2^{-m}\}\bigr)<2^{-m}.
\]
The limsup of these exceptional sets has capacity zero, since for every $N$
it is contained in their union over $m\geq N$, whose capacity is at most
$2^{1-N}$.  Outside the limsup, the desired convergence holds.
\end{proof}

\begin{remark}
\label{rem:capacity-subsequence-sequential}
The proposition is deliberately sequential.  A convergent net in capacity
need not contain a useful sequence unless the directed set or the capacity
space has an additional countability property.  Later compactness theorems
will provide exactly the required reductions.
\end{remark}

Both hypotheses of \Cref{prop:capacity-subsequence} are used: countable
subadditivity bounds the capacity of the union $\bigcup_{m\geq N}$, and the
sequential indexing makes that union countable.  Neither is available for a
net, which is why the statement is not upgraded.  The reduction from a net to
a sequence, when it is possible at all, is performed before this proposition
is applied, by \Cref{lem:choquet-selection} or by an explicit cofinal
sequence.

\section{Concave transforms and Duistermaat--Heckman distributions}
\label{sec:dh-preview}

Let $\Delta\subset\mathbf R^n$ be a compact convex body and let
$\lambda_\Delta$ denote Lebesgue measure restricted to $\Delta$.  We do not
normalize it: its total mass is $\operatorname{vol}(\Delta)$.  Let
\[
               G:\Delta\longrightarrow\mathbf R\cup\{-\infty\}
\]
be a proper upper semicontinuous concave function.  Its effective domain
$\{G> -\infty\}$ is convex.

We first record the facts about concave functions on a convex body that make
the construction legitimate.  Here \emph{proper} means, as in
\Cref{def:appD-proper-closed} applied to $-G$, that $G$ is finite at some
point and never equal to $+\infty$.

\begin{lemma}
\label{lem:concave-on-body}
Let $\Delta\subset\mathbf R^n$ be a compact convex body and let
$G:\Delta\to\mathbf R\cup\{-\infty\}$ be proper, concave, and upper
semicontinuous.  Then:
\begin{enumerate}
\item every superlevel set $\{G\geq\tau\}$ and every strict superlevel set
  $\{G>\tau\}$ is convex, the former being closed;
\item $G$ is bounded above and attains its maximum;
\item $G$ is Borel measurable, and it is continuous on the relative interior
  of its effective domain $D=\{G>-\infty\}$;
\item $G$ is Lebesgue integrable on $\Delta$ if and only if
  $\int_\Delta\max(-G,0)\,\rd x<\infty$, and in that case
  $\lambda_\Delta(\{G=-\infty\})=0$.
\end{enumerate}
\par
\end{lemma}

\begin{proof}
\emph{(1).}  If $G(x),G(y)\geq\tau$ and $s\in[0,1]$, concavity gives
$G((1-s)x+sy)\geq(1-s)G(x)+sG(y)\geq\tau$, and the same computation with
strict inequalities gives the second assertion.  The set $\{G\geq\tau\}$ is
closed because $\{G<\tau\}$ is open, $G$ being usc.

\emph{(2).}  This is \Cref{cor:usc-attains-max}, since $\Delta$ is compact and
nonempty and $G$ takes values in $\mathbf R\cup\{-\infty\}$.

\emph{(3).}  Upper semicontinuity makes every $\{G<\tau\}$ open, hence Borel,
so $G$ is Borel measurable.  For continuity on the relative interior of $D$,
extend $G$ by $-\infty$ outside $\Delta$, so that $-G$ becomes a proper convex
function on $\mathbf R^n$ with effective domain $D$.  A convex function is
continuous relative to any relatively open convex subset of its effective
domain, by \cite[Theorem~10.1]{Roc70}; applying this to $\RelInt D$ gives the
assertion.

\emph{(4).}  Write $G=G^+-G^-$ with $G^+=\max(G,0)$ and $G^-=\max(-G,0)$.
By (2), $G^+$ is bounded, hence integrable on the bounded set $\Delta$, so
integrability of $G$ is equivalent to integrability of $G^-$.  If $G^-$ is
integrable then $\{G^-=+\infty\}=\{G=-\infty\}$ is Lebesgue-null.
\end{proof}

\begin{remark}
\label{rem:concave-boundary}
Part (3) says nothing about the boundary of $D$, and nothing can be said: on
$\Delta=[0,1]$ the function $G(x)=\sqrt{x}$ is concave, continuous, and
finite, while $G(x)=\log x$ for $x>0$ with $G(0)=-\infty$ is concave and usc
with $D=(0,1]$, and it is unbounded below at the boundary point $0$.  Both are
integrable, so both define Duistermaat--Heckman distributions; only the second
has a nontrivial tail near $-\infty$, namely
$\DHm_{[0,1],\log}\bigl((-\infty,\tau]\bigr)=\rme^{\tau}$ for $\tau\leq0$.  In
\cref{chap:transcendental-okounkov-dh}, where $G$ arises as a concave
transform, exactly this boundary behaviour encodes the singularities of a
potential.
\end{remark}

\begin{definition}
\label{def:dh-distribution}
Assume that $G$ is finite almost everywhere, and let
$N=\{G=-\infty\}$, a null set.  Choose any finite-valued Borel function
$\widetilde G$ that agrees with $G$ on $\Delta\setminus N$.  The
\emph{Duistermaat--Heckman distribution} of the pair $(\Delta,G)$ is
\[
                        \operatorname{DH}_{\Delta,G}
                        :=\widetilde G_*\lambda_\Delta.
\]
It is a positive Radon measure on $\mathbf R$ of total mass
$\operatorname{vol}(\Delta)$ and is independent of the chosen values on
$N$, because $N$ is null.  Indeed it is a finite Borel measure, and every
finite Borel measure on $\mathbf R$ is Radon; continuity of $\widetilde G$ is
neither assumed nor needed.
\end{definition}

This normalization is designed for later Okounkov-body constructions.  If
the total Monge--Amp\`ere mass is $(L^n)=n!\operatorname{vol}(\Delta)$, the
corresponding unnormalized geometric distribution is
$n!\operatorname{DH}_{\Delta,G}$.

The same pushforward formula defines $\operatorname{DH}_{\Delta,G}$ for any
finite-valued Borel function $G$ on $\Delta$.  Concavity is needed for the
tail geometry, not for the next two convergence statements.

\begin{proposition}\label{prop:dh-tail}
For every $\tau\in\mathbf R$,
\[
 \operatorname{DH}_{\Delta,G}((\tau,+\infty))
 =\operatorname{vol}\{x\in\Delta:G(x)>\tau\}.
\]
If
$V(\tau)=\operatorname{vol}\{G>\tau\}$, then $V$ is decreasing and
right-continuous, and $\operatorname{DH}_{\Delta,G}$ is its negative
Stieltjes derivative.  Atoms record jumps of $V$.
\end{proposition}

\begin{proof}
The first identity is the definition of pushforward.  The second is the
standard correspondence between a finite measure on $\mathbf R$ and its tail
function.

In detail, choose the finite Borel representative $\widetilde G$ from the
definition.  Then
$\widetilde G_*\lambda_\Delta\bigl((\tau,+\infty)\bigr)
=\lambda_\Delta\bigl(\widetilde G^{-1}(\tau,+\infty)\bigr)
=\vol\{x\in\Delta:G(x)>\tau\}$: the two inverse images can differ only on
the null set $\{G=-\infty\}$, and the last set is Borel by
\Cref{lem:concave-on-body}\,(3).  The function $V$ is
nonincreasing because the sets $\{G>\tau\}$ decrease with $\tau$.  It is
right-continuous because
$\{G>\tau\}=\bigcup_{m\geq1}\{G>\tau+1/m\}$ is an increasing union, so
$\lambda_\Delta$ evaluated on this increasing union gives
$V(\tau)=\lim_mV(\tau+1/m)$.  Consequently $V$ is the tail function of the finite Borel measure
$\DHm_{\Delta,G}$ on $\mathbf R$, which is therefore determined by $V$: the
measure of an interval $(\tau_0,\tau_1]$ is $V(\tau_0)-V(\tau_1)$, and these
intervals generate the Borel $\sigma$-algebra.  Left limits are given by
\[
 V(\tau^-)=\lim_{m}V(\tau-1/m)=\vol\{G\geq\tau\},
\]
because $\{G\geq\tau\}=\bigcap_m\{G>\tau-1/m\}$ is a decreasing intersection
of sets of finite measure.  Hence
\[
 \DHm_{\Delta,G}(\{\tau\})=V(\tau^-)-V(\tau)
 =\vol\{G\geq\tau\}-\vol\{G>\tau\}=\vol\{G=\tau\},
\]
which is the assertion about atoms.
\end{proof}

\begin{corollary}
\label{cor:dh-affine}
Let $c\in\mathbf R$ and $\lambda>0$.  Writing $T_c(t)=t+c$ and
$m_\lambda(t)=\lambda t$ for the corresponding maps $\mathbf R\to\mathbf R$,
\[
 \DHm_{\Delta,G+c}=(T_c)_*\DHm_{\Delta,G},
 \qquad
 \DHm_{\Delta,\lambda G}=(m_\lambda)_*\DHm_{\Delta,G}.
\]
Both operations preserve the total mass $\vol(\Delta)$.
\end{corollary}

\begin{proof}
For a Borel set $B\subseteq\mathbf R$,
\[
 \DHm_{\Delta,G+c}(B)=\vol\{x:G(x)+c\in B\}
 =\vol\{x:G(x)\in B-c\}=\DHm_{\Delta,G}(T_c^{-1}B),
\]
and similarly for $\lambda G$.  The mass statement is the case
$B=\mathbf R$, using that $T_c$ and $m_\lambda$ are bijections of
$\mathbf R$.
\end{proof}

The Duistermaat--Heckman distribution is narrowly continuous under uniform
convergence of the concave transform.  This is the only continuity statement
used in \cref{chap:transcendental-okounkov-dh}, and it is elementary; the
delicate continuity statements there concern the dependence of $G$ itself on
the geometric data.

\begin{proposition}
\label{prop:dh-continuity}
Let $\Delta$ be a compact convex body and let $G_j,G:\Delta\to\mathbf R$ be
bounded Borel functions with $G_j\to G$ uniformly on $\Delta$.  Then
$\DHm_{\Delta,G_j}\to\DHm_{\Delta,G}$ narrowly on $\mathbf R$, that is,
\[
 \int_{\mathbf R}\chi\,\rd\DHm_{\Delta,G_j}
 \longrightarrow
 \int_{\mathbf R}\chi\,\rd\DHm_{\Delta,G}
\]
for every bounded continuous $\chi:\mathbf R\to\mathbf R$.
\end{proposition}

\begin{proof}
All the measures are supported in a fixed compact interval $I\subseteq\mathbf
R$: indeed $\lVert G_j\rVert_\infty\leq\lVert G\rVert_\infty+1$ for $j$ large.
Let $\chi$ be bounded and continuous; on the compact interval $I$ it is
uniformly continuous, so given $\varepsilon>0$ there is $\delta>0$ with
$|\chi(s)-\chi(t)|<\varepsilon$ whenever $s,t\in I$ and $|s-t|<\delta$.  By
\Cref{lem:dh-moments}, applied to the bounded Borel function $\chi$,
\[
 \Bigl|\int_{\mathbf R}\chi\,\rd\DHm_{\Delta,G_j}
 -\int_{\mathbf R}\chi\,\rd\DHm_{\Delta,G}\Bigr|
 =\Bigl|\int_\Delta\bigl(\chi(G_j(x))-\chi(G(x))\bigr)\rd x\Bigr|
 \leq\varepsilon\vol(\Delta)
\]
as soon as $\lVert G_j-G\rVert_\infty<\delta$.
\end{proof}

\begin{corollary}
\label{cor:dh-continuity-ae}
Let $G_j,G:\Delta\to\mathbf R$ be Borel functions with $G_j\to G$ Lebesgue
almost everywhere on $\Delta$.  Then
$\DHm_{\Delta,G_j}\to\DHm_{\Delta,G}$ narrowly on $\mathbf R$.
\end{corollary}

\begin{proof}
Let $\chi:\mathbf R\to\mathbf R$ be bounded and continuous.  Then
$\chi\circ G_j\to\chi\circ G$ almost everywhere on $\Delta$, and
$|\chi\circ G_j|\leq\lVert\chi\rVert_\infty$, which is integrable on the
bounded set $\Delta$.  Dominated convergence and \Cref{lem:dh-moments} give
\[
 \int_{\mathbf R}\chi\,\rd\DHm_{\Delta,G_j}
 =\int_\Delta\chi(G_j)\,\rd x
 \longrightarrow\int_\Delta\chi(G)\,\rd x
 =\int_{\mathbf R}\chi\,\rd\DHm_{\Delta,G}.
\]
Since $\chi$ was arbitrary, this is narrow convergence.
\end{proof}

\begin{remark}[Narrow convergence does not give convergence of moments]
\label{rem:dh-moments-not-continuous}
Narrow convergence of Duistermaat--Heckman distributions controls only bounded
continuous test functions, and the unbounded test function $\chi(t)=t$, whose
integral is the barycenter, is genuinely not continuous for it.  On
$\Delta=[0,1]$ put
\[
 G_j(x)=\min\bigl(0,\,-j^2+j^4x\bigr),\qquad j\geq1 ,
\]
a minimum of two affine functions, hence concave, with $G_j(0)=-j^2$ and
$G_j\equiv0$ on $[j^{-2},1]$.  Then $G_j\to0$ pointwise on $(0,1]$, so
$\DHm_{[0,1],G_j}\to\delta_0$ narrowly by \Cref{cor:dh-continuity-ae}.
Explicitly, $\DHm_{[0,1],G_j}$ is the sum of an atom of mass $1-j^{-2}$ at
$0$ and the uniform measure of density $j^{-4}$ on $[-j^2,0]$.  Its barycenter
is
\[
 \int_{\mathbf R}t\,\rd\DHm_{[0,1],G_j}(t)
 =\int_0^1G_j(x)\,\rd x=-\tfrac12
\]
for every $j$, whereas the barycenter of $\delta_0$ is $0$.  Mass escapes to
$-\infty$ while carrying a fixed first moment.  This is the reason that
\cref{chap:transcendental-okounkov-dh} states energy and entropy statements
with explicit uniform integrability or monotonicity hypotheses, and never
deduces convergence of a moment from narrow convergence alone.
\end{remark}

\begin{proposition}\label{prop:dh-brunn-minkowski}
Assume $\Delta$ has nonempty interior.  On the interval where it is positive,
the function
\[
               \tau\longmapsto
               \operatorname{vol}\{G\geq\tau\}^{1/n}
\]
is concave.
\end{proposition}

\begin{proof}
Concavity of $G$ gives
\[
 (1-t)\{G\geq\tau_0\}+t\{G\geq\tau_1\}
 \subseteq
 \{G\geq(1-t)\tau_0+t\tau_1\}.
\]
Apply the Brunn--Minkowski inequality.

The inclusion is immediate: if $G(x_0)\geq\tau_0$ and $G(x_1)\geq\tau_1$ then
concavity gives
$G((1-t)x_0+tx_1)\geq(1-t)G(x_0)+tG(x_1)\geq(1-t)\tau_0+t\tau_1$.  The sets
$\{G\geq\tau_i\}$ are compact and convex by
\Cref{lem:concave-on-body}\,(1), and they are nonempty on the interval where
the volume is positive.  \Cref{cor:appD-brunn-minkowski}, applied to the
convex bodies $(1-t)\{G\geq\tau_0\}$ and $t\{G\geq\tau_1\}$ and combined with
monotonicity of volume under the displayed inclusion, gives
\[
 \vol\{G\geq(1-t)\tau_0+t\tau_1\}^{1/n}
 \geq(1-t)\vol\{G\geq\tau_0\}^{1/n}+t\,\vol\{G\geq\tau_1\}^{1/n},
\]
using the homogeneity $\vol(sK)=s^n\vol(K)$.  This is the asserted concavity.
\end{proof}

\begin{remark}
\label{rem:dh-log-concavity}
\Cref{prop:dh-brunn-minkowski} is the source of every log-concavity statement
about Duistermaat--Heckman measures in this book.  It says that the closed-tail
function $W(\tau)=\vol\{G\geq\tau\}$ has a concave $n$-th root on the interval
where $W>0$.  The strict tail $V(\tau)=\vol\{G>\tau\}$ is its right-continuous
version and differs from $W$ precisely at the atoms of $\DHm_{\Delta,G}$; no
concavity assertion is made across those jumps.  In dimension $n=1$,
concavity of $W$ implies that the absolutely continuous density
$-W'$ is nondecreasing wherever it is defined; an atom at a terminal plateau
is recorded separately by the jump from $W$ to $V$.  The corresponding
statement in \cref{chap:bj-test-configurations} for Duistermaat--Heckman
measures of test configurations is proved by the same Brunn--Minkowski
mechanism after passing to Okounkov bodies, not by a separate argument.
\end{remark}

\begin{example}[Affine transform]
\label{ex:dh-affine-transform}
For $\Delta=[a,b]$ and $G(x)=x$, one has
$\operatorname{DH}_{\Delta,G}=\mathbf 1_{[a,b]}\,\rd t$.  If instead
$\Delta=[0,1]$ and $G(x)=\min\{x,1-x\}$, then
\[
 \operatorname{DH}_{\Delta,G}=2\mathbf 1_{[0,1/2]}\,\rd t.
\]
Indeed, the superlevel interval has length $1-2\tau$ for
$0\leq\tau\leq1/2$.
\end{example}

\begin{example}[A piecewise affine transform on a simplex]
\label{ex:dh-simplex}
Let $\Delta\subset\mathbf R^2$ be the standard simplex with vertices
$(0,0)$, $(1,0)$, $(0,1)$, so that $\vol(\Delta)=1/2$, and let
\[
 G(x_1,x_2)=\min\{x_1,\ x_2,\ c\},\qquad 0<c<1/2 .
\]
As a minimum of affine functions, $G$ is concave and continuous on $\Delta$,
with values in $[0,c]$.  We compute $\DHm_{\Delta,G}$ from
\Cref{prop:dh-tail}.

For $0\leq\tau<c$ the superlevel set is
\[
 \{G>\tau\}=\{x\in\Delta:x_1>\tau,\ x_2>\tau\},
\]
a triangle similar to $\Delta$ with ratio $1-2\tau$, so
\[
 V(\tau)=\vol\{G>\tau\}=\tfrac12(1-2\tau)^2 .
\]
For $\tau\geq c$ we have $\{G>\tau\}=\varnothing$ and $V(\tau)=0$.  Hence $V$
is smooth on $[0,c)$ with $V'(\tau)=-2(1-2\tau)$, and it jumps at $\tau=c$
from $V(c^-)=\tfrac12(1-2c)^2$ to $0$.  By \Cref{prop:dh-tail},
\[
 \DHm_{\Delta,G}
 =2(1-2\tau)\,\mathbf 1_{[0,c)}(\tau)\,\rd\tau
 +\tfrac12(1-2c)^2\,\delta_c .
\]
The total mass is as it must be:
\[
 \int_0^c2(1-2\tau)\,\rd\tau+\tfrac12(1-2c)^2
 =(2c-2c^2)+\bigl(\tfrac12-2c+2c^2\bigr)=\tfrac12=\vol(\Delta).
\]

Three features are worth recording, because all three recur in
\cref{chap:bj-test-configurations} and
\cref{chap:transcendental-okounkov-dh}.  The atom at $\tau=c$ has mass
$\vol\{G=c\}$, in accordance with the atom formula in the proof of
\Cref{prop:dh-tail}; it appears exactly because $G$ is constant on a set of
positive measure, which is what a truncation $\min(\cdot,c)$ produces.  The
density $2(1-2\tau)$ is decreasing, in accordance with
\Cref{prop:dh-brunn-minkowski} in dimension $n=2$: indeed
$V(\tau)^{1/2}=2^{-1/2}(1-2\tau)$ is affine, the extremal case of concavity.
Finally, replacing $G$ by $G+c'$ translates the whole picture, by
\Cref{cor:dh-affine}, and does not change the shape of the distribution.
\end{example}

\begin{lemma}\label{lem:dh-moments}
Let $\widetilde G$ be any finite Borel representative used in the definition
of $\operatorname{DH}_{\Delta,G}$.  For every Borel function $\chi$ that is
integrable on the right-hand side,
\[
 \int_{\mathbf R}\chi(t)\,\rd \operatorname{DH}_{\Delta,G}(t)
 =\int_\Delta\chi(\widetilde G(x))\,\rd x.
\]
In particular, if $G$ is integrable, the barycenter is
\[
 \frac{1}{\operatorname{vol}(\Delta)}
 \int_\Delta G(x)\,\rd x.
\]
\end{lemma}

\begin{proof}
This is the change-of-variables identity defining a pushforward measure.
Changing $\widetilde G$ on the null set $\{G=-\infty\}$ changes neither
side.  When $G$ is integrable, its extended Lebesgue integral is likewise
unchanged by the chosen finite values on that null set.
\end{proof}

\section{Extended-real arithmetic}
\label{sec:extended-arithmetic}

Potentials may equal $-\infty$, energies may equal $-\infty$, and dual
energies may equal $+\infty$.  We use the following ledger throughout.

The ledger below is not a list of conventions chosen for convenience.  Each
entry is forced, in the sense that the alternative destroys a statement the
book needs; and each entry is used, in the sense that some later proof would
be false without it.  After stating the ledger we go through it entry by
entry, giving the justification and an explicit example of what goes wrong if
the entry is dropped.  The reader who wants the short version may consult
\Cref{tab:extended-real-ledger}.

\begin{enumerate}
\item The order on $\overline{\mathbf R}$ is
  $-\infty<a<+\infty$ for $a\in\mathbf R$.
\item For $a\in\mathbf R$,
  \[
  a+(+\infty)=+\infty,\qquad a+(-\infty)=-\infty.
  \]
  We also set $(+\infty)+(+\infty)=+\infty$ and
  $(-\infty)+(-\infty)=-\infty$; the mixed sum
  $(+\infty)+(-\infty)$ is undefined.
\item For $c>0$, set $c(+\infty)=+\infty$ and
  $c(-\infty)=-\infty$.  The expressions $0(\pm\infty)$ are undefined unless
  an explicit recession convention has been introduced.
\item A difference $a-b$ means $a+(-b)$ only when this does not form
  $(+\infty)+(-\infty)$.  In particular, $(-\infty)-(-\infty)$ and
  $(+\infty)-(+\infty)$ are undefined.
\item We set
  \[
             \sup\varnothing=-\infty,\qquad
             \inf\varnothing=+\infty.
  \]
\item For a net $(a_i)$ in $\overline{\mathbf R}$,
  \[
  \limsup_i a_i=\inf_{i_0}\sup_{i\geq i_0}a_i,\qquad
  \liminf_i a_i=\sup_{i_0}\inf_{i\geq i_0}a_i.
  \]
\item The integral of an extended-real function is used only when at least
  one of its positive and negative parts has finite integral.  We never
  assign a value to $+\infty-\infty$.
\end{enumerate}

\begin{table}
\caption{The extended-real ledger, with the failure that each entry prevents.}
\label{tab:extended-real-ledger}
\begin{tabular}{@{}lll@{}}
\toprule
Operation & Convention & Failure without it\\
\midrule
order & $-\infty<a<+\infty$ &
  no monotone limits, no $\sup$/$\inf$\\
$a+(\pm\infty)$, $a$ finite & $\pm\infty$ &
  $\varphi+c$ undefined on the polar locus\\
$(+\infty)+(-\infty)$ & undefined &
  $0=\infty$ by associativity\\
$c\cdot(\pm\infty)$, $c>0$ & $\pm\infty$ &
  $\lambda\varphi$ undefined for a psh $\varphi$\\
$0\cdot(\pm\infty)$ & undefined outside integrals &
  discontinuity of $\lambda\mapsto\lambda\varphi$ hidden\\
$0\cdot(+\infty)$ inside an integral & $0$ &
  null sets could carry mass\\
$a-b$ & only if not $(+\infty)+(-\infty)$ &
  $\varphi-\psi=0$ on a polar set gives $\varphi=\psi$\\
$\sup\varnothing$, $\inf\varnothing$ & $-\infty$, $+\infty$ &
  envelopes over empty families undefined\\
$\limsup$, $\liminf$ of a net & tail $\inf\sup$, $\sup\inf$ &
  no upper limit for a divergent net\\
$\int f\,\rd\mu$ & one of $f^\pm$ integrable &
  $+\infty-\infty$ in an energy identity\\
\bottomrule
\end{tabular}
\end{table}

\subsection{The ledger entry by entry}
\label{subsec:ledger-commentary}

\emph{(1) The order.}  The order makes $\Rbar$ a complete lattice: every
subset has a supremum and an infimum.  This is what guarantees that
$\sup_i\varphi_i$, $\inf_a f_a$, $f^*$, $\limsup$, and $\liminf$ are defined
without any hypothesis, and it is used in
\Cref{lem:monotone-net-converges} and in every appearance of
\eqref{eq:usc-regularization}.  The order topology on $\Rbar$ is compact and
metrizable; a homeomorphism onto $[-1,1]$ is given by
$t\mapsto t/(1+|t|)$ extended by $\pm1$, which shows in particular that a
monotone net of extended reals converges and that
\Cref{prop:appE-real-cluster} applies to arbitrary $\Rbar$-valued nets.
Without the order, none of the envelope calculus of
\cref{sec:usc-envelopes} exists.

\emph{(2) Addition.}  Setting $a+(-\infty)=-\infty$ for finite $a$ is forced
by monotonicity: adding a constant must be an order isomorphism of $\Rbar$,
and it is the reason that
$\varphi\mapsto\varphi+c$ is well defined on potentials that are $-\infty$
somewhere, which is used already in the normalization of energies.  The mixed
sum must be left undefined.  Suppose one declared
$(+\infty)+(-\infty)=s$ for some $s\in\Rbar$.  Associativity would then give
\[
 s=\bigl((+\infty)+(-\infty)\bigr)+(-\infty)
  =(+\infty)+\bigl((-\infty)+(-\infty)\bigr)=(+\infty)+(-\infty)=s,
\]
which is harmless, but the same computation with $1+(+\infty)=+\infty$ gives
\[
 s=(+\infty)+(-\infty)=\bigl(1+(+\infty)\bigr)+(-\infty)
  =1+\bigl((+\infty)+(-\infty)\bigr)=1+s ,
\]
which forces $s=\pm\infty$ and already destroys cancellation, on which
\Cref{rem:forbidden-cancellations} depends.  Associativity by itself does
not exclude a value: the ``inf-addition'' convention
$(+\infty)\oplus(-\infty)=-\infty$ is commutative, monotone and associative,
and so is its mirror.  What no choice survives is compatibility with limits:
for every $c\in\mathbf R$ the sequences $r_n=n\to+\infty$ and
$s_n=-n+c\to-\infty$ have $r_n+s_n\to c$, so no single value of $s$ is the
limit of the sums.  This is \Cref{lem:extended-real-safe}\,(5), and it is
why the mixed sum is left undefined rather than fixed by convention.

\emph{(3) Positive scaling and the case $c=0$.}  For $c>0$, multiplication by
$c$ is an order isomorphism of $\Rbar$, which is what makes
\Cref{lem:usc-star-calculus}\,(2) and
\Cref{prop:appE-safe-identities} true.  The product $0\cdot(-\infty)$ is left
undefined because the two natural values disagree.  Concretely, let
$\varphi(z)=\log|z|$ on the unit disc.  For $\lambda>0$ the function
$\lambda\varphi$ is subharmonic with $\lambda\varphi(0)=-\infty$, so the
pointwise limit as $\lambda\downarrow0$ is
\[
 \lim_{\lambda\downarrow0}\lambda\varphi(z)=
 \begin{cases}0,&z\neq0,\\ -\infty,&z=0,\end{cases}
\]
whereas the regularized limit is the constant $0$, and the ``value at
$\lambda=0$'' suggested by the usc-envelope topology is again $0$.
Declaring $0\cdot(-\infty)=0$ agrees with the regularized limit but not with
the pointwise limit at the origin; declaring it to be $-\infty$ agrees with
the pointwise limit but not with the regularized one.  Thus no pointwise
convention is continuous for both limiting procedures.  We leave the product
undefined and, in \cref{chap:trivial-analytification}, state the topology and
the regularization used by each scaling limit, exactly as in
\Cref{ex:sup-not-usc-log}.

\emph{(4) Differences.}  A difference is formed only when it is not of the
shape $(+\infty)+(-\infty)$.  The failure this prevents is the most common
error in the subject.  Let $\varphi(z)=\log|z|$ and $\psi(z)=2\log|z|$ on the
unit disc.  These are distinct subharmonic functions, yet at the origin
``$\varphi-\psi$'' would be $(-\infty)-(-\infty)$, and a reader who set that
equal to $0$ would conclude $\varphi=\psi$ at the origin from
$\varphi-\psi=0$ there.  The same mechanism invalidates the naive proof that
two potentials with the same Monge--Amp\`ere measure differ by a constant, and
it is the reason \cref{chap:bfj-energy} normalizes potentials before
subtracting.  See \Cref{rem:forbidden-cancellations} for the three
cancellations that fail.

\emph{(5) Empty suprema and infima.}  The conventions
$\sup\varnothing=-\infty$ and $\inf\varnothing=+\infty$ are forced by
monotonicity of $\sup$ and $\inf$ under inclusion of index sets, and they are
used whenever an envelope is formed over a family that may be empty:
\[
 P(\varphi)=\Bigl(\sup\bigl\{\psi\in\PSH(X,\theta):\psi\leq\varphi\bigr\}\Bigr)^*
\]
is then defined even when no psh function lies below $\varphi$, with value
$-\infty$, which is the correct answer.  Without the convention, the obstacle
envelope of \cref{chap:bfj-envelopes} would be undefined precisely in the
degenerate cases that its monotonicity statements have to cover, and an
induction on a decreasing family of nonempty competitor sets could not pass to
the limit.

\emph{(6) Upper and lower limits of nets.}  These are defined by the tail
formulas, which make sense for every net because of entry (1).  They are
\emph{not} defined as limits of subnets, although
\Cref{prop:appE-real-cluster} identifies them with the extreme subnet limits
when the net is bounded.  The tail definition is what makes
\eqref{eq:net-usc-test} an equivalence rather than an implication and what
makes $\limsup$ monotone under passage to subnets.

\emph{(7) Integrals.}  An integral is formed only when at least one of
$\int f^+$ and $\int f^-$ is finite.  Inside integrals, and only there, we use
$0\cdot(+\infty)=0$, justified by monotone truncation: if $\mu(A)=0$ then
\[
 \int_A f\,\rd\mu=\lim_{M\to\infty}\int_A\min(f,M)\,\rd\mu=0
\]
for $f\geq0$, by ordinary monotone convergence, since each truncated integral
vanishes.
Without this convention a set of measure zero on which a potential equals
$-\infty$ would contribute an undefined term to
$\int\varphi\,\rd\mu$, and the finite-energy classes of
\cref{chap:bfj-energy} could not be defined by finiteness of that integral.
The convention is confined to integrals: outside them,
$0\cdot(\pm\infty)$ remains undefined by entry (3), and the difference
between the two regimes is exactly the content of the conventions sheet.

\begin{lemma}
\label{lem:extended-real-safe}
Let $(r_i)_{i\in I}$ be a nonempty family in $\Rbar$, let $c\in\mathbf R$, and
let $\lambda>0$.  Then, with all expressions defined by the ledger:
\begin{enumerate}
\item $\sup_i(r_i+c)=\sup_ir_i+c$ and $\inf_i(r_i+c)=\inf_ir_i+c$;
\item $\sup_i(\lambda r_i)=\lambda\sup_ir_i$ and
  $\inf_i(\lambda r_i)=\lambda\inf_ir_i$;
\item $\sup_i(-r_i)=-\inf_ir_i$;
\item for nets $(r_a)$, $(s_a)$ in $\mathbf R\cup\{-\infty\}$ indexed by the
  same directed set, and provided the right-hand side is not of the form
  $(+\infty)+(-\infty)$,
\[
 \limsup_a(r_a+s_a)\leq\limsup_ar_a+\limsup_as_a ,
\]
  and the inequality may be strict;
\item if $r_a\to r$ in $\Rbar$ with $r$ finite and $s_a\to s$ in $\Rbar$, then
  $r_a+s_a\to r+s$, the sum being defined.
\end{enumerate}
\par
\end{lemma}

\begin{proof}
\emph{(1)--(3)} are \Cref{prop:appE-safe-identities}: translation by a finite
constant, multiplication by $\lambda>0$, and negation are order isomorphisms
of $\Rbar$, hence commute with suprema and infima, negation exchanging the
two.

\emph{(4).}  Fix $a_0$ and write $R_{a_0}=\sup_{b\geq a_0}r_b$ and
$S_{a_0}=\sup_{b\geq a_0}s_b$, both in $\mathbf R\cup\{\pm\infty\}$.  If
either is $-\infty$, then the corresponding net is identically $-\infty$ on
that tail, hence $\sup_{a\geq a_0}(r_a+s_a)=-\infty$ and the inequality is
trivial.  Otherwise $R_{a_0}+S_{a_0}$ is defined and
$r_a+s_a\leq R_{a_0}+S_{a_0}$ for $a\geq a_0$, so
$\sup_{a\geq a_0}(r_a+s_a)\leq R_{a_0}+S_{a_0}$.  Both $R_{a_0}$ and
$S_{a_0}$ decrease as $a_0$ increases, with infima
$R=\limsup_ar_a$ and $S=\limsup_as_a$.  If $R+S=+\infty$, the desired
inequality is automatic.  If $R,S$ are finite, then for every
$\varepsilon>0$ both tail suprema are eventually below $R+\varepsilon$ and
$S+\varepsilon$, respectively, so
$\limsup_a(r_a+s_a)\leq R+S+2\varepsilon$; let
$\varepsilon\downarrow0$.  It remains to consider $R+S=-\infty$.  After
interchanging the two nets, we may assume $R=-\infty$ and $S<+\infty$.
Choose a tail on which $S_{a_0}\leq B$ for some $B\in\mathbf R$.  Given
$T\in\mathbf R$, refine it until $R_{a_0}<T-B$; on that tail
$r_a+s_a<T$.  Since $T$ was arbitrary, the left-hand limsup is $-\infty$.
These cases exhaust all defined values of $R+S$.  For strictness take
$r_n=(-1)^n$ and $s_n=(-1)^{n+1}$ in $\mathbf R$: the left side is $0$ and
the right side is $2$.

\emph{(5).}  If $s$ is finite this is continuity of addition on
$\mathbf R\times\mathbf R$ together with the observation that both nets are
eventually finite.  If $s=+\infty$, then for every $M$ the net $s_a$ is
eventually $>M+1-r$ and $r_a$ is eventually $>r-1$, so $r_a+s_a>M$ eventually,
whence $r_a+s_a\to+\infty=r+s$.  The case $s=-\infty$ is symmetric.  Note that
finiteness of $r$ is essential: entry (2) of the ledger leaves $r+s$ undefined
otherwise, and the conclusion fails, as $r_n=n$, $s_n=-n$ shows.
\end{proof}

\begin{lemma}\label{lem:safe-supremum}
Let $(a_i)$ and $(b_j)$ be families in
$\mathbf R\cup\{-\infty\}$, both bounded above by real numbers.  Then
\[
 \sup_{i,j}(a_i+b_j)=\sup_i a_i+\sup_j b_j.
\]
If $c>0$, then $\sup_i(ca_i)=c\sup_i a_i$.  The same statements hold with
infima in $\mathbf R\cup\{+\infty\}$ and families bounded below.
\end{lemma}

\begin{proof}
The inequality $\leq$ is immediate.  Approximate each finite supremum from
below; if one supremum is $-\infty$, every member of that family is
$-\infty$ and both sides equal $-\infty$.  Positive homogeneity is equally
direct.  The infimum statements follow by changing signs.
\end{proof}

\begin{remark}[Three forbidden cancellations]
\label{rem:forbidden-cancellations}
The following implications are invalid without finiteness hypotheses:
\[
 a+c=b+c\Longrightarrow a=b,
 \qquad
 a-b=0\Longrightarrow a=b,
 \qquad
 \lim(a_i-b_i)=\lim a_i-\lim b_i.
\]
For example, the first implication carries no information when
$c=+\infty$.  Later energy proofs will normalize potentials before
subtracting quantities that may be infinite.
\end{remark}

Explicit witnesses for the three failures, in the form in which they arise
later, are the following.  For the first, take $a=0$, $b=1$, $c=+\infty$: then
$a+c=b+c=+\infty$ while $a\neq b$.  In context, this is the observation that
two potentials with the same infinite energy cannot be compared by cancelling
the energy.  For the second, take $a=b=-\infty$ and read $a-b$ as $0$: this is
the failure discussed in entry (4) of the ledger, and it is the reason that
$\varphi\leq\psi$ and $\psi\leq\varphi$ are proved separately rather than
through $\varphi-\psi=0$.  For the third, take $a_i=b_i=-i$ in $\mathbf R$:
then $a_i-b_i=0$ for every $i$, so the left side is $0$, while the right side
is $(-\infty)-(-\infty)$ and is undefined.  In context, this is the reason
that slopes of energies along rays are defined by a single limit of
differences rather than as a difference of two limits, as in
\cref{chap:rays-filtrations}.

\begin{lemma}
\label{lem:integration-ledger}
Let $(X,\mu)$ be a finite measure space and let $f,g:X\to\Rbar$ be
measurable.
\begin{enumerate}
\item If $f$ is bounded above by a real constant, then
  $\int_Xf\,\rd\mu$ is well defined in $[-\infty,+\infty)$, and
  $\int_X(f+c)\,\rd\mu=\int_Xf\,\rd\mu+c\,\mu(X)$ for every $c\in\mathbf R$.
\item If $f$ is bounded above, then
  $\int_X\max(f,-M)\,\rd\mu\searrow\int_Xf\,\rd\mu$ as $M\to+\infty$, both
  sides possibly being $-\infty$.
\item If $\mu(A)=0$, then $\int_Af\,\rd\mu=0$ for every measurable $f$,
  including $f\equiv-\infty$ on $A$.
\item If $f\leq g$ and both integrals are defined, then
  $\int_Xf\,\rd\mu\leq\int_Xg\,\rd\mu$.  If moreover
  $\int_Xg\,\rd\mu$ is finite, then $\int_X(g-f)\,\rd\mu$ is defined and
  equals $\int_Xg\,\rd\mu-\int_Xf\,\rd\mu$.
\end{enumerate}
\par
\end{lemma}

\begin{proof}
\emph{(1).}  Boundedness above by a constant $C$ gives
$f^+\leq\max(C,0)$, so that
\[
 \int_Xf^+\,\rd\mu\leq\max(C,0)\,\mu(X)<\infty,
\]
and the integral is defined in $[-\infty,+\infty)$ by ledger entry (7).  The
translation identity holds for bounded $f$ by ordinary linearity; in general
apply it to $\max(f,-M)$ and let $M\to\infty$, using (2).

\emph{(2).}  Write $f_M=\max(f,-M)$.  The functions $f_M^-$ increase to
$f^-$ as $M\to\infty$, so that $\int_Xf_M^-\,\rd\mu$ increases to
$\int_Xf^-\,\rd\mu$ in $[0,+\infty]$ by monotone convergence, while
$f_M^+=f^+$ has finite integral.  Subtracting a
finite quantity from an increasing one gives the assertion.

\emph{(3).}  Write $f=f^+-f^-$.  For $h\geq0$ measurable,
$\int_Ah\,\rd\mu=\sup\{\int_A s\,\rd\mu: 0\leq s\leq h\ \text{simple}\}$, and
every simple $s$ has $\int_As\,\rd\mu=0$ because $\mu(A)=0$; monotone
truncation therefore gives $\int_Ah\,\rd\mu=0$ even where $h=+\infty$.  This
is exactly the convention $0\cdot(+\infty)=0$ of ledger entry (7), and it
applies to $f^+$ and $f^-$ separately.

\emph{(4).}  Monotonicity holds because $f^+\leq g^+$ and $f^-\geq g^-$,
by monotonicity of the integral of nonnegative functions.  For the difference,
finiteness of $\int_Xg\,\rd\mu$ forces both $\int_Xg^+\,\rd\mu$ and
$\int_Xg^-\,\rd\mu$ to be finite, so $g$ is finite $\mu$-almost everywhere and
$\int_Xf^+\,\rd\mu\leq\int_Xg^+\,\rd\mu$ is finite.  Hence $g-f\geq0$ is
defined outside a $\mu$-null set, on which (3) applies, and it is measurable.
If $\int_Xf\,\rd\mu=-\infty$, then $f=-\infty$ on a set of positive measure or
$\int f^-=+\infty$, and in either case $\int_X(g-f)\,\rd\mu=+\infty$, which
agrees with $\int_Xg\,\rd\mu-\int_Xf\,\rd\mu=+\infty$.  If
$\int_Xf\,\rd\mu$ is finite, then $f$ is finite almost everywhere and
$g=f+(g-f)$ almost everywhere, so additivity of the integral for functions
with finite integral gives the identity.
\end{proof}

\begin{remark}[Where the ledger is invoked]
\label{rem:ledger-uses}
The entries are used as follows.  Entry (1) underlies every monotone limit
in this chapter.  Entry (2) makes the normalization
$\varphi\mapsto\varphi+c$ of \cref{chap:bfj-energy} legitimate on potentials
with polar singularities.  Entry (3) is what allows a potential to be scaled
in \cref{chap:trivial-analytification} and \cref{chap:bj-homogeneity-bdivisors}
without a case distinction, once the scaling limit is stated for regularized
limits.  Entry (4) is invoked at every comparison of two singular potentials
and is the reason for the relative formulation
$\varphi\leq\psi+C$ rather than $\varphi-\psi\leq C$.  Entry (5) makes
envelopes total operations.  Entry (6) is the definition behind
\Cref{prop:net-semicontinuity}.  Entry (7), with
\Cref{lem:integration-ledger}, is what makes the finite-energy classes
$\cE^1$ well defined.
\end{remark}

\section{Dictionary and source note}

The recurring correspondences established in this chapter are
\[
\begin{array}{ccl}
\text{compactness} &\longleftrightarrow&
   \text{convergent subnets},\\
\text{supremum of potentials} &\longmapsto&
   \text{usc-regularized envelope},\\
\text{second countability} &\longmapsto&
   \text{Choquet sequence selection},\\
\text{continuous test functions} &\longleftrightarrow&
   \text{weak convergence of Radon measures},\\
\text{concave energy} &\longmapsto&
   \text{convex dual energy of measures},\\
\text{concave transform on a body} &\longmapsto&
   \text{Duistermaat--Heckman distribution}.
\end{array}
\]

The list continues with the correspondences added in this second layer of the
chapter:
\[
\begin{array}{ccl}
\text{upper semicontinuity} &\longleftrightarrow&
   \limsup_a f(x_a)\leq f(x)\ \text{along nets},\\[1mm]
\text{decreasing net of usc functions} &\longmapsto&
   \text{usc infimum, monotone integral},\\[1mm]
\text{positive linear functional on }C(X) &\longleftrightarrow&
   \text{Radon measure},\\[1mm]
\text{concave function on }\mathbf R^n &\longleftrightarrow&
   \text{convex conjugate on }(\mathbf R^n)^\vee,\\[1mm]
\text{free energy} &\longleftrightarrow&
   \text{relative entropy},\\[1mm]
\text{countably subadditive capacity} &\longmapsto&
   \sigma\text{-ideal of polar sets},\\[1mm]
\text{atom of }\DHm_{\Delta,G} &\longleftrightarrow&
   \text{level set of }G\text{ of positive volume}.
\end{array}
\]

\subsection*{Checklist of statements exported to later chapters}

For convenience of reference, the statements of this chapter that are quoted
elsewhere are the following.  Compactness and convergence:
\Cref{prop:net-criterion}, \Cref{prop:net-semicontinuity},
\Cref{cor:usc-attains-max}, \Cref{lem:monotone-net-converges}.  Dini:
\Cref{prop:dini-nets} and \Cref{prop:dini-usc}, the latter in part (2) more
often than in part (1).  Regularization:
\Cref{prop:usc-majorant}, \Cref{prop:usc-limsup-description},
\Cref{lem:usc-star-calculus}, \Cref{lem:usc-inf}.  Selection:
\Cref{lem:choquet-selection} and \Cref{prop:choquet-regularizing}.  Measures:
\Cref{thm:riesz-interface}, \Cref{cor:riesz-uniqueness},
\Cref{lem:pushforward-radon}, \Cref{prop:measure-compactness},
\Cref{prop:portmanteau}, \Cref{lem:joint-weak-uniform},
\Cref{cor:usc-net-monotone-integral}.  Duality:
\Cref{lem:dual-mass}, \eqref{eq:fenchel-energy},
\Cref{prop:concave-biconjugation}.  Capacity:
\Cref{prop:polar-sigma-ideal}, \Cref{prop:quasicontinuity-calculus},
\Cref{prop:capacity-subsequence}.  Distributions:
\Cref{prop:dh-tail}, \Cref{prop:dh-brunn-minkowski},
\Cref{prop:dh-continuity}, \Cref{lem:dh-moments}.  Arithmetic:
\Cref{lem:extended-real-safe}, \Cref{lem:safe-supremum},
\Cref{lem:integration-ledger}.

\subsection*{Source note}

The net and upper-envelope conventions are the ones used in the BFJ and BJ
global theories; see \cite{BFJ16,BJ22}.  The unnormalized
Duistermaat--Heckman convention agrees with the intersection-theoretic
normalization in \cite{BHJ17}.  Concrete capacities and their
quasicontinuity theorems will be proved only after the corresponding
Monge--Amp\`ere operators have been constructed.

The point-set material is classical.  Moore--Smith convergence, subnets, and
the equivalence between compactness and the existence of convergent subnets go
back to the 1920s and 1930s; the treatment we follow, including the tail
formulation of subnets, is Willard's \cite[Chapters~4--5]{Wil70}, with
Kelley \cite{Kel75} as a second reference and as the source of the Tychonoff
theorem in the form used in \Cref{ex:net-no-subsequence} and
\Cref{prop:measure-compactness}.  The terminological variation surrounding
Willard and Moore--Smith subnets is recorded in \cref{app:nets}; throughout
this book the defining condition is the order-preserving tail condition
\eqref{eq:appE-willard}.  The compact space
$\{0,1\}^{[0,1]}$ of \Cref{ex:net-no-subsequence} and the one-point
compactification of an uncountable discrete set used in
\Cref{ex:decreasing-net-not-sequential} and
\Cref{ex:choquet-fails-uncountable} are the two standard test spaces for
sequential arguments, and both are found in the topology literature; we have
given complete proofs of the properties we use rather than citing them,
because the details are exactly what the later chapters must check.

Upper semicontinuous regularization and the selection lemma of
\Cref{lem:choquet-selection} are, in the complex pluripotential setting, part
of the standard toolkit; a textbook account with the same conventions is
\cite[Chapter~1 and Chapter~4]{GZ17}, and the classical source for the
regularization theorem in potential theory is Brelot's study of subharmonic
functions at a singular point \cite{Bre49}.  The name ``Choquet lemma'' is
standard in the pluripotential literature.  Our
\Cref{prop:choquet-regularizing} isolates the countability hypothesis; the
form in which it is used in the non-Archimedean setting, where the ambient
space may fail to be metrizable, is the reason for that isolation, and the
corresponding discussion in the trivially valued theory appears in
\cite{BJ22}.

The Riesz representation theorem is quoted from
\cite[Theorems~2.14 and~6.19]{Rud87}; \cite[Chapter~7]{Fol99} gives the same
material with the regularity hypotheses displayed as in
\eqref{eq:radon-regularity}, and is the reference we recommend for the reader
who wants the construction rather than the statement.  Weak convergence of
measures, the portmanteau inequalities, and the compactness of
$\mathcal P(X)$ are classical; the proof given here is the direct
Tychonoff--Riesz argument, deliberately avoiding Banach--Alaoglu, so that no
functional-analytic apparatus beyond \Cref{thm:riesz-interface} is needed
anywhere in this book.

Finite-dimensional convex duality follows \cite[Chapters~10, 12, and~23]{Roc70},
and we have kept the sign conventions of \cref{app:convexity} so that the
concave conjugate \eqref{eq:concave-conjugate} matches
\eqref{eq:appD-conjugate} after a change of sign.  The Brunn--Minkowski
inequality used in \Cref{prop:dh-brunn-minkowski} is
\Cref{cor:appD-brunn-minkowski}, proved in \cref{app:convexity} from the
Alexandrov--Fenchel inequality, for which the reference is \cite{Sch14}.  The
free-energy/relative-entropy pair of \Cref{ex:legendre-entropy-pair} is the
finite model of the Gibbs variational principle; in the pluripotential
literature it is the finite-dimensional shadow of the energy/entropy duality
that governs the variational approach to Monge--Amp\`ere equations, developed
in \cite{BBGZ13} in the complex setting and imitated in the non-Archimedean
setting in \cite{BFJ15,BJ22}.  We have chosen the finite model because it can
be computed completely, and because every hypothesis of the general theory is
visible in it as a nondegeneracy of the finite computation.

Capacities in the sense used here go back to Choquet; the axioms we state are
the ones needed for the Monge--Amp\`ere capacities of \cref{part:bfj} and
\cref{part:bj}, whose complex ancestor is the Bedford--Taylor capacity of
\cite{BT82}, with the quasicontinuity theory developed there and in
\cite{GZ17}.  We have kept the treatment axiomatic because the several
capacities constructed later satisfy the axioms for different reasons, and
because the statements that fail without countable subadditivity, as in
\Cref{ex:capacity-not-subadditive}, are exactly the ones most often used
without comment.

Duistermaat--Heckman measures entered this circle of ideas through
\cite{BHJ17}, whose normalization we follow; the concave-transform description
used here, and the corresponding Okounkov-body constructions, are treated in
\cref{chap:transcendental-okounkov-dh}.  The interaction between atoms,
truncations, and log concavity illustrated in
\Cref{ex:dh-simplex,rem:dh-moments-not-continuous} is the elementary shadow of
the phenomena studied there.

Finally, a word on the extended-real ledger.  Conventions for arithmetic on
$\Rbar$ are usually left implicit, and the resulting errors are not exotic:
the failures collected in \Cref{rem:forbidden-cancellations} and in the
commentary of \cref{subsec:ledger-commentary} all occur in the literature.  We
have therefore fixed the conventions in the frontmatter, repeated them here in
full, and given for each entry an explicit example of what its omission
destroys.  The general-purpose safe identities are collected once more, in
their most general form, in \Cref{prop:appE-safe-identities}.
